%!TEX TS-program = xelatex
%!TEX encoding = UTF-8 Unicode
% Compile with XeLaTeX (Tectonic): the Lean listings in Appendices B-C are Unicode and use fontspec.
\documentclass[11pt,a4paper]{article}
\usepackage[margin=1.05in]{geometry}
\usepackage{amsmath,amssymb,amsthm}
\usepackage{booktabs}
\usepackage{microtype}
\usepackage{fontspec}
\usepackage{fancyvrb}
\usepackage{newunicodechar}
%% The Lean listings in Appendices B and C contain Unicode mathematical symbols
%% (ℝ, ∑, ⊗, →, δ, μ, …).  They are set in \leanmono, a monospaced font that
%% carries them; on systems whose available monospaced font lacks some of them,
%% those glyphs alone are taken from a fallback font.  The rest of the document
%% is unaffected.
\newif\ifleanfallback\leanfallbackfalse
\IfFontExistsTF{DejaVu Sans Mono}{\newfontfamily\leanmono{DejaVu Sans Mono}}{%
  \IfFontExistsTF{Consolas}{%
    \newfontfamily\leanmono{Consolas}%
    \IfFontExistsTF{Cambria Math}{%
      \newfontfamily\leanfallbackfont{Cambria Math}\leanfallbacktrue}{}%
  }{\let\leanmono\ttfamily}%
}
\ifleanfallback
\newunicodechar{𝔼}{{\leanfallbackfont\symbol{120124}}}
\newunicodechar{ₖ}{{\leanfallbackfont\symbol{8342}}}
\newunicodechar{ℕ}{{\leanfallbackfont\symbol{8469}}}
\newunicodechar{ℝ}{{\leanfallbackfont\symbol{8477}}}
\newunicodechar{↦}{{\leanfallbackfont\symbol{8614}}}
\newunicodechar{∀}{{\leanfallbackfont\symbol{8704}}}
\newunicodechar{∃}{{\leanfallbackfont\symbol{8707}}}
\newunicodechar{∈}{{\leanfallbackfont\symbol{8712}}}
\newunicodechar{∧}{{\leanfallbackfont\symbol{8743}}}
\newunicodechar{∼}{{\leanfallbackfont\symbol{8764}}}
\newunicodechar{⊗}{{\leanfallbackfont\symbol{8855}}}
\newunicodechar{⊢}{{\leanfallbackfont\symbol{8866}}}
\newunicodechar{⟨}{{\leanfallbackfont\symbol{10216}}}
\newunicodechar{⟩}{{\leanfallbackfont\symbol{10217}}}
\fi
\usepackage[hidelinks]{hyperref}
\emergencystretch=1.5em
\newtheorem{theorem}{Theorem}[section]
\newtheorem{lemma}[theorem]{Lemma}
\newtheorem{proposition}[theorem]{Proposition}
\newtheorem{corollary}[theorem]{Corollary}
\theoremstyle{definition}
\newtheorem{definition}[theorem]{Definition}
\newtheorem{hypothesis}[theorem]{Hypothesis}
\theoremstyle{remark}
\newtheorem{remark}[theorem]{Remark}
\newcommand{\E}{\mathbb{E}}
\newcommand{\Prob}{\mathbb{P}}
\newcommand{\R}{\mathbb{R}}
\newcommand{\F}{\mathcal{F}}
\newcommand{\cC}{\mathcal{C}}
\newcommand{\cM}{\mathcal{M}}
\newcommand{\ent}{h}
\newcommand{\iid}{\mathrm{iid}}
\newcommand{\id}{\mathrm{id}}
\DeclareMathOperator{\Cov}{Cov}
\DeclareMathOperator{\Var}{Var}
\DeclareMathOperator{\supp}{supp}

\title{The ceiling of the single-letter entropy method for the union-closed sets conjecture, and a protocol that reaches it}
\author{Andrew Moffat\thanks{Moffat Studio, United Kingdom. Email: \texttt{andrewmoffat.business@gmail.com}.}}
\date{8 September 2026; revised 10 September 2026 (Appendices B--C added)}

\begin{document}
\maketitle

\begin{abstract}
Gilmer's entropy method for Frankl's union-closed sets conjecture, in the form developed by Sawin, Yu, Cambie and Liu, certifies a constant $c$ (some element lies in at least $c|\F|$ members of every union-closed family $\F$) through a single-letter inequality. The inequality involves a \emph{protocol}, a memoryless rule that couples the two conditional bits, and a \emph{class} of admissible joint laws of the two prefix-conditional probabilities. We prove that every such certificate whose classes contain product laws is bounded by $1-\ent(1/\sqrt2)/\sqrt2=0.383099\ldots$, and that every certificate which uses the i.i.d.\ protocol and whose other classes admit \emph{component hiding} (a property of every class in the literature) is bounded by $c^{**}=0.382885260\ldots$, the joint solution of two explicit adversarial laws. We then construct a conditionally-i.i.d.\ protocol whose kernel has the optimal diagonal and show, by a computer-assisted argument of the type used for the previous record $0.382709$ of Liu, with the small-entropy regime handled by a new analytic bound, that it certifies $c=0.38284$, and, under the analogous hypothesis at each weight, every $c<c^{**}$. The improvement is small, $+1.8\cdot10^{-4}$, and the second theorem shows that it is the last one available to this family of arguments. The two ceiling theorems (Theorems~\ref{thm:ceil} and~\ref{thm:refined}) and Lemma~\ref{lem:classes}, with the exception of its maximal-correlation case, have been formalised and machine-checked in Lean~4 with Mathlib; the formal development is included as ancillary material and is described in Appendix~\ref{app:lean}.
\end{abstract}

\noindent\textit{2020 Mathematics Subject Classification.} Primary 05D05; Secondary 94A17, 60E15.\\
\textit{Keywords.} Union-closed sets conjecture, Frankl's conjecture, entropy method, coupling, maximal correlation.

\section{Introduction}\label{sec:intro}
A family $\F$ of finite sets is \emph{union-closed} if $A\cup B\in\F$ whenever $A,B\in\F$. Frankl's conjecture \cite{Frankl,BruhnSchaudt} asserts that if $\F$ is a finite union-closed family other than $\{\emptyset\}$, then some element belongs to at least half of the members of $\F$. Write $c(\F)$ for the largest fraction $|\{A\in\F:i\in A\}|/|\F|$ over elements $i$, and let $c_0$ be the largest constant such that $c(\F)\ge c_0$ for all such $\F$; the conjecture is $c_0=\tfrac12$, and the family $\{\emptyset,\{1\}\}$ shows $c_0\le\tfrac12$.

In November 2022 Gilmer \cite{Gilmer} proved $c_0\ge0.01$ by an entropy argument, the first constant lower bound. Within days Alweiss, Huang and Sellke \cite{AHS}, Chase and Lovett \cite{CL}, Sawin \cite{Sawin} and Pebody \cite{Pebody} sharpened the argument to $c_0\ge\psi:=(3-\sqrt5)/2=0.381966\ldots$; Chase and Lovett showed that $\psi$ is the correct constant for the \emph{approximate} version of the conjecture, in which $A\cup B\in\F$ is only required for most pairs. Sawin proposed coupling the two random sets in Gilmer's argument; Yu \cite{Yu} and Cambie \cite{Cambie} carried this out, and Cambie proved that Sawin's coupling gives exactly $0.3823455\ldots$. Liu \cite{Liu} introduced \emph{conditionally i.i.d.} couplings and obtained $c_0\ge0.382709$, conditional on two numerically verified hypotheses. Table~\ref{tab:history} summarises the record.

\begin{table}[h]\centering\small
\begin{tabular}{llp{5.6cm}}\toprule
lower bound on $c_0$ & reference & mechanism\\\midrule
$0.01$ & Gilmer \cite{Gilmer} & entropy of the union of two independent samples\\
$(3-\sqrt5)/2=0.381966$ & \cite{AHS,CL,Sawin,Pebody} & sharp form of the same inequality; optimal for approximate closure\\
$0.3823455$ & Sawin \cite{Sawin}, Yu \cite{Yu}, Cambie \cite{Cambie} & anti-correlated coupling; exact value by Cambie\\
$0.382709$ & Liu \cite{Liu} & conditionally-i.i.d.\ coupling (computer-assisted)\\
$0.38284$, and $c<0.382885$ & this paper & optimal-diagonal conditionally-i.i.d.\ protocol (computer-assisted)\\\bottomrule
\end{tabular}
\caption{Lower bounds on the union-closed constant.}\label{tab:history}
\end{table}

All of these arguments have the same shape, made explicit by Liu \cite[Definition~1, Proposition~3]{Liu}: a protocol, a class of admissible couplings, and a single-letter inequality (see Section~\ref{sec:framework}). Our first contribution is to determine the supremum of the constants that this framework can certify.

\begin{theorem}[informal; Theorems~\ref{thm:ceil} and~\ref{thm:refined}]\label{thm:A}
Every single-letter certificate whose classes contain product laws certifies at most $1-\ent(1/\sqrt2)/\sqrt2=0.383099\ldots$. Every single-letter certificate that uses the i.i.d.\ protocol and whose other classes admit component hiding certifies at most $c^{**}=0.382885260\ldots$. All classes used in \cite{Gilmer,AHS,CL,Sawin,Pebody,Yu,Cambie,Liu} contain product laws and admit hiding.
\end{theorem}

Both ceiling theorems, together with the parts of Lemma~\ref{lem:classes} that concern the
i.i.d.\ class, the class of all couplings and the class of mixtures of products, have been
formalised and machine-checked in Lean~4 with Mathlib, on a finite model of the framework of
Section~\ref{sec:framework}; the maximal-correlation case of Lemma~\ref{lem:classes} is not
formalised. The formal development proves, for every certificate satisfying the hypotheses of
Theorem~\ref{thm:refined} with $c\le\tfrac12$, that $c\le0.3829$: the exact fixed point
$c^{**}$ is a real number evaluated numerically, whereas the rational bound $0.3829$ is
proved outright, and since $0.3829<c_{\mathrm{ceil}}=0.3830993\ldots$ it already separates the
two ceilings formally. The development is included as ancillary material; it is described in
Appendix~\ref{app:lean} and its complete source is reproduced in Appendix~\ref{app:src}.

Our second contribution is a protocol that reaches this ceiling. Write $x=\Prob(\text{bit}=0)$ for a conditional zero-probability.

\begin{theorem}[informal; Theorem~\ref{thm:main}]\label{thm:B}
Let $f(x)=\min(x,1-x)$ for $x\le\frac12$, $f(x)=\sqrt{\frac12-x^2}$ for $\frac12\le x\le\frac1{\sqrt2}$ and $f(x)=0$ otherwise, and let $\Pi^{\id}$ be the conditionally-i.i.d.\ protocol with $\Prob(\text{both bits }0)=xy+f(x)f(y)$. Under two numerically verified hypotheses of the same kind as those in \cite{Liu}, the mixture of the i.i.d.\ protocol with weight $w=0.810222$ and $\Pi^{\id}$ with weight $1-w$ certifies $c_0\ge0.38284$; letting $w$ increase to the value forced by Theorem~\ref{thm:A}, the same computation, under the analogous hypothesis at each weight, certifies every $c<c^{**}$.
\end{theorem}

The proofs of the ceiling theorems are short and elementary: two explicit families of laws of the conditional probabilities defeat every protocol. The first, a two-point law, is the one Gilmer's argument has always been fighting. The second, which we call \emph{component hiding}, is new: once the two random sets are coupled in a way that the adversary may rearrange, an atom of tiny mass that carries all the entropy can be kept away from the atom on which the element is surely absent, so that only the i.i.d.\ term, whose class is a singleton, is guaranteed the cross pairs that carry entropy. This forces the i.i.d.\ weight to be at least $1/(2(1-c))\approx0.81$ and is what makes the remaining improvement small.

\paragraph{Organisation.} Section~\ref{sec:framework} recalls Gilmer's argument and Liu's framework, with a self-contained proof of the certificate. Section~\ref{sec:ceiling} proves the ceiling theorems. Section~\ref{sec:protocol} constructs the protocol and computes its value against the two adversaries. Section~\ref{sec:card} gives the cardinality reduction (a generalisation of Liu's Theorem~12) and the small-entropy bound. Section~\ref{sec:numerics} describes the computation and states the main numerical theorem. Section~\ref{sec:discussion} discusses what would be needed to go further. Appendix~\ref{app:const} records the constants.

\section{Gilmer's argument and the single-letter framework}\label{sec:framework}
Throughout, $\ent(p)=-p\log_2p-(1-p)\log_2(1-p)$ is the binary entropy, $H$ is Shannon entropy in bits, and $[n]=\{1,\dots,n\}$. For $\F\subseteq2^{[n]}$ we identify a set with its indicator string in $\{0,1\}^n$.

\subsection{Gilmer's argument}
Let $\F\subseteq2^{[n]}$ be union-closed, $\F\ne\{\emptyset\}$, and let $A=(A_1,\dots,A_n)$ be uniform on $\F$, so $H(A)=\log_2|\F|$. Let $B$ be a second random set with the same law, coupled with $A$ in some way. Since $A\cup B\in\F$,
\begin{equation}\label{eq:union}
H(A\cup B)\le\log_2|\F|=H(A).
\end{equation}
For a prefix $a\in\{0,1\}^{i-1}$ with $\Prob(A_{<i}=a)>0$ write
\[x_i(a):=\Prob(A_i=0\mid A_{<i}=a)\in[0,1]\]
for the conditional \emph{zero}-probability. By the chain rule and the fact that $(A\cup B)_{<i}$ is a function of $(A_{<i},B_{<i})$,
\begin{equation}\label{eq:chain}
\begin{aligned}
H(A\cup B)&=\sum_{i=1}^nH\big((A\cup B)_i\mid(A\cup B)_{<i}\big)\;\ge\;\sum_{i=1}^nH\big((A\cup B)_i\mid A_{<i},B_{<i}\big),\\
H(A)&=\sum_{i=1}^n\E\,\ent\big(x_i(A_{<i})\big).
\end{aligned}
\end{equation}
Conditionally on $(A_{<i},B_{<i})=(a,b)$ the bit $(A\cup B)_i$ is Bernoulli with $\Prob((A\cup B)_i=0\mid a,b)=\Prob(A_i=0,B_i=0\mid a,b)$, so the $i$-th summand on the left is $\E\,\ent(\Prob(A_i=B_i=0\mid A_{<i},B_{<i}))$. Gilmer takes $B$ independent of $A$, so that $\Prob(A_i=B_i=0\mid a,b)=x_i(a)x_i(b)$, and proves an inequality of the form $\E\,\ent(XY)\ge\E\,\ent(X)$ for i.i.d.\ $X,Y$ with $\E X$ large; the sharp version \cite{AHS,CL,Sawin,Pebody} is $\ent(xy)\ge\frac{\varphi}2\,(x\ent(y)+y\ent(x))$ with $\varphi=(1+\sqrt5)/2$, tight at $x=y=1/\varphi$, which gives $\E\,\ent(XY)\ge\varphi\,\E[X]\,\E\,\ent(X)\ge\E\,\ent(X)$ whenever $\E X\ge1/\varphi$, i.e.\ whenever every element has frequency at most $1-1/\varphi=\psi$.

\subsection{Protocols, classes and the certificate}
\begin{definition}[protocol, class; \cite{Liu}]\label{def:protocol}
A \emph{protocol} $\Pi$ assigns to every $(x,y)\in[0,1]^2$ a probability law $\Pi_{x,y}$ on $\{0,1\}^2$ whose two marginals have zero-probabilities $x$ and $y$. Given $\F$, the protocol generates $(A,B)$ coordinate by coordinate: after $(A_{<i},B_{<i})=(a,b)$ has been generated, $(A_i,B_i)$ is drawn from $\Pi_{x_i(a),x_i(b)}$. A \emph{class} for $\Pi$ is a map $\mu\mapsto\cC_\Pi(\mu)$ from laws on $[0,1]$ to sets of couplings of $\mu$ with itself (laws on $[0,1]^2$ with both marginals $\mu$) such that, for every union-closed $\F$ and every $i$, the joint law of $\big(x_i(A_{<i}),x_i(B_{<i})\big)$ belongs to $\cC_\Pi(\mu)$, where $\mu$ is the law of $x_i(A_{<i})$.
\end{definition}
The marginal laws of $A$ and $B$ under any protocol are both uniform on $\F$: summing the product $\prod_i\Pi_{x_i(a_{<i}),x_i(b_{<i})}(a_i,b_i)$ over $b_n,b_{n-1},\dots$ in turn leaves $\prod_i\Prob(A_i=a_i\mid A_{<i}=a_{<i})$.

\begin{proposition}[the certificate; \cite{Liu}, Proposition~3]\label{prop:cert}
Let $\Pi^{(1)},\dots,\Pi^{(K)}$ be protocols with classes $\cC_1,\dots,\cC_K$, let $w_1,\dots,w_K\ge0$ sum to $1$, and let $c\in(0,\tfrac12]$. Suppose there is $C>1$ such that
\begin{equation}\label{eq:cert}
\sum_{k=1}^Kw_k\inf_{P\in\cC_k(\mu)}\E_{(X,Y)\sim P}\Big[\ent\big(\Pi^{(k)}_{X,Y}(0,0)\big)\Big]\;\ge\;C\,\E_\mu\big[\ent(X)\big]
\end{equation}
for every law $\mu$ on $[0,1]$ with $\E_\mu[X]\ge1-c$. Then $c(\F)\ge c$ for every finite union-closed $\F\ne\{\emptyset\}$.
\end{proposition}
\begin{proof}
Suppose $c(\F)<c$. For each $i$ let $\mu_i$ be the law of $x_i(A_{<i})$; then $\E_{\mu_i}[X]=\Prob(A_i=0)>1-c$. Generate $(A,B^{(k)})$ with protocol $\Pi^{(k)}$. By \eqref{eq:chain}, the definition of a class, and \eqref{eq:cert},
\[\sum_kw_kH(A\cup B^{(k)})\ge\sum_i\sum_kw_k\inf_{P\in\cC_k(\mu_i)}\E_P\,\ent\big(\Pi^{(k)}_{X,Y}(0,0)\big)\ge C\sum_i\E_{\mu_i}\ent(X)=C\,H(A).\]
By \eqref{eq:union} the left side is at most $H(A)$, so $H(A)=0$ and $|\F|=1$, i.e.\ $\F=\{S\}$; then $c(\F)=1\ge c$ unless $S=\emptyset$, which is excluded.
\end{proof}

The four protocol/class pairs used so far are the following; here and below $\mu\otimes\mu$ denotes the product law.
\begin{enumerate}
\item \emph{I.i.d.} $\Pi^{\iid}_{x,y}(0,0)=xy$, class $\{\mu\otimes\mu\}$ \cite{Gilmer}.
\item \emph{Sawin's maximum-entropy protocol.} $\Pi_{x,y}(0,0)=1-\max\{1-x,\,1-y,\,\min(2-x-y,\tfrac12)\}$, the point of the Fr\'echet interval nearest $\frac12$ when $x,y\ge\frac12$ and comonotone otherwise; class: all couplings of $\mu$ with itself \cite{Sawin,Yu,Cambie}.
\item \emph{Yu's maximal-correlation protocols.} $\Pi_{x,y}$ with maximal correlation at most $\rho$; class: couplings of $\mu$ with itself of maximal correlation at most $\rho$ (tensorisation of maximal correlation) \cite{Yu}.
\item \emph{Liu's conditionally-i.i.d.\ protocols.} The two bits are generated conditionally independently given a shared uniform variable $U$: $\Prob(A_i=0\mid U=u)=r_x(u)$ with $r_x:[0,1]\to[0,1]$, $\int_0^1r_x=x$, and likewise for $B_i$ with $r_y$. Then
\begin{equation}\label{eq:kernel}
\Pi_{x,y}(0,0)=\int_0^1r_x(u)r_y(u)\,du=xy+K(x,y),\qquad K(x,y)=\Cov(r_x,r_y),
\end{equation}
where $K$ is a positive semidefinite kernel; class $\cC_3(\mu)$: the laws $\int\nu\otimes\nu\,d\pi(\nu)$, mixtures of product laws over a probability measure $\pi$ on laws $\nu$ on $[0,1]$, whose marginal $\int\nu\,d\pi(\nu)$ equals $\mu$ \cite{Liu}. Liu's Theorem~9 shows that for every conditionally-i.i.d.\ protocol the infimum in \eqref{eq:cert} over $\cC_3(\mu)$ is attained at a mixture of two products,
\begin{equation}\label{eq:twoprod}
P=(1-q)\,P_0\otimes P_0+q\,P_1\otimes P_1,\qquad\mu=(1-q)P_0+qP_1,\qquad q\in[0,1].
\end{equation}
Liu's own protocol is the rank-one kernel $K(x,y)=x(1-x)\,y(1-y)$, realised by $r_x(u)=x+x(1-x)\,\mathrm{sgn}(u-\tfrac12)$.
\end{enumerate}

\section{Two adversaries and the ceiling}\label{sec:ceiling}
Two facts about protocols are used repeatedly. If $x=0$ then $A_i=1$ surely, so the union bit is surely $1$ whatever the protocol: pairs containing the atom $0$ carry no entropy. On any pair $(x,y)$, $\Pi_{x,y}(0,0)$ lies in the Fr\'echet interval $[\max(0,x+y-1),\min(x,y)]$ and $\ent(\Pi_{x,y}(0,0))\le1$.

\begin{theorem}[product ceiling]\label{thm:ceil}
Suppose every class $\cC_k(\mu)$ in \eqref{eq:cert} contains the product law $\mu\otimes\mu$ for every $\mu$. Then \eqref{eq:cert} fails for every $c>c_{\mathrm{ceil}}:=1-\ent(1/\sqrt2)/\sqrt2=0.3830993\ldots$
\end{theorem}
\begin{proof}
Let $x^*=1/\sqrt2$, so $\ent(x^{*2})=\ent(\tfrac12)=1$, and let $\mu=p\,\delta_{x^*}+(1-p)\,\delta_0$ with $p\in(0,1)$. Then $\E_\mu[X]=px^*$ and $\E_\mu[\ent(X)]=p\,\ent(x^*)$. For each $k$ the infimum in \eqref{eq:cert} is at most the value at $\mu\otimes\mu$, which is at most $p^2\cdot1$: the pair $(x^*,x^*)$ has mass $p^2$ and entropy at most $1$, and every other pair contains the atom $0$. Hence \eqref{eq:cert} implies $p^2\ge Cp\,\ent(x^*)>p\,\ent(x^*)$, i.e.\ $p>\ent(1/\sqrt2)$. Choose $p$ slightly smaller than $\ent(1/\sqrt2)$; the law $\mu$ violates \eqref{eq:cert}, and its mean $p/\sqrt2$ is slightly smaller than $\ent(1/\sqrt2)/\sqrt2=1-c_{\mathrm{ceil}}$, hence at least $1-c$ for every $c>c_{\mathrm{ceil}}$ once $p$ is close enough.
\end{proof}

The i.i.d.\ protocol itself has $\ent(x^2)=1$ at $x=1/\sqrt2$; its constant $\psi$ is smaller than $c_{\mathrm{ceil}}$ only because for $x<1/\sqrt2$ it under-uses the diagonal, $\ent(x^2)<1$. Every improvement since 2022 consists in recovering part of that deficit while controlling what an enlarged class allows the adversary to do. The second adversary bounds how much can be recovered.

\begin{definition}[component hiding]\label{def:hiding}
For $q,y\in(0,1)$ and $\delta\in(0,\tfrac12)$ let
\begin{equation}\label{eq:mudelta}
\mu_{q,y,\delta}:=q\,\delta_1+(1-q)\big[(1-\delta)\,\delta_0+\delta\,\delta_y\big].
\end{equation}
A class $\cC$ \emph{admits hiding} if for all $q,y\in(0,1)$ and all sufficiently small $\delta$, $\cC(\mu_{q,y,\delta})$ contains a coupling $P$ with
\[P(X=y,\,Y=1)=P(X=1,\,Y=y)=0\qquad\text{and}\qquad P(X=y,\,Y=y)\le\delta^2 .\]
\end{definition}

\begin{lemma}\label{lem:classes}
The classes of the i.i.d.\ protocol, of Sawin's protocol, of Yu's protocols for every $\rho>0$, and of every conditionally-i.i.d.\ protocol contain the product law. All but the first admit hiding.
\end{lemma}
\begin{proof}
Products: immediate for the singleton, for all couplings, and for mixtures of products; the product law has maximal correlation $0\le\rho$. Hiding for the class of all couplings is trivial (take the coupling of the maximal-correlation case below, which has $P(X=y,Y=y)=0$). For mixtures of products take $P=q\,\delta_1\otimes\delta_1+(1-q)\,P_0\otimes P_0$ with $P_0=(1-\delta)\delta_0+\delta\delta_y$; its marginal is $\mu_{q,y,\delta}$, $P(X=y,Y=1)=0$ and $P(X=y,Y=y)=(1-q)\delta^2\le\delta^2$. For the maximal-correlation class, write $p_1=q$, $p_0=(1-q)(1-\delta)$, $p_y=(1-q)\delta$ and define the coupling $P$ on $\{1,0,y\}^2$ by $P(y,0)=P(0,y)=p_y$, $P(y,y)=P(y,1)=P(1,y)=0$, and on $\{1,0\}^2$ by $P(a,b)=r_ar_b/(r_1+r_0)$ with $r_1=p_1$, $r_0=p_0-p_y$. Both marginals are $(p_1,p_0,p_y)$. The maximal correlation of $P$ is the second largest singular value of the matrix $M_{ab}=P(a,b)/\sqrt{p_ap_b}$. As $\delta\to0$, $r_0\to p_0$ and the $\{1,0\}$ block of $M$ converges to the rank-one matrix $\sqrt{p_ap_b}$ of the product law, while the entries of the row and column of $y$ are $0$ or $\sqrt{p_y/p_0}\to0$. Singular values depend continuously on the entries, so the second singular value tends to $0$ and is below $\rho$ for all small $\delta$.
\end{proof}

\begin{theorem}[refined ceiling]\label{thm:refined}
Suppose \eqref{eq:cert} uses the i.i.d.\ protocol with its singleton class $\{\mu\otimes\mu\}$ and weight $w\in[0,1]$, and that every other class contains the product law and admits hiding. Then
\begin{align}
2w(1-c)&\ge1,\label{eq:H}\\
c&\le1-\frac{x\,\ent(x)}{w\,\ent(x^2)+1-w}\qquad\text{for every }x\in[\tfrac12,\tfrac1{\sqrt2}]\text{ with }\ent(x)\le w\,\ent(x^2)+1-w.\label{eq:D}
\end{align}
Consequently $c\le c^{**}$, where $c^{**}=0.382885260\ldots$ is the unique solution of $c=G(c)$,
\begin{equation}\label{eq:G}
G(c):=\min_{x\in X_c}\Big[1-\frac{x\,\ent(x)}{1-(1-\ent(x^2))/(2(1-c))}\Big],\qquad X_c:=\Big\{x\in[\tfrac12,\tfrac1{\sqrt2}]:\ent(x)\le1-\tfrac{1-\ent(x^2)}{2(1-c)}\Big\}.
\end{equation}
\end{theorem}
\begin{proof}
\emph{Inequality \eqref{eq:H}.} Fix $y\in(0,1)$, put $q=1-c$ and $\mu_\delta=\mu_{q,y,\delta}$. Then $\E_{\mu_\delta}[X]=q+(1-q)\delta y\ge1-c$ and $\E_{\mu_\delta}[\ent(X)]=(1-q)\delta\,\ent(y)$. Consider any coupling of $\mu_\delta$ with itself and any protocol. Pairs containing the atom $0$ carry no entropy. On the pair $(1,1)$ both bits are surely $0$, so the union bit is surely $0$ and carries no entropy. On the pairs $(y,1)$ and $(1,y)$ one bit is surely $0$, so the union bit equals the other bit and has entropy exactly $\ent(y)$, for every protocol. The pair $(y,y)$ carries entropy at most $1$; its mass is $(1-q)^2\delta^2$ under the product law and at most $\delta^2$ under a hiding coupling, by definition.

For the i.i.d.\ term, whose class is the singleton $\{\mu_\delta\otimes\mu_\delta\}$, the value is exactly $2q(1-q)\delta\,\ent(y)+(1-q)^2\delta^2\,\ent(y^2)$. For every other protocol the infimum over its class is at most its value at a hiding coupling, which is at most $\delta^2$. Hence
\[\text{LHS of \eqref{eq:cert}}\le w\big[2q(1-q)\delta\,\ent(y)+(1-q)^2\delta^2\ent(y^2)\big]+(1-w)\,\delta^2 .\]
Dividing by $\E_{\mu_\delta}\ent(X)=(1-q)\delta\,\ent(y)$ and letting $\delta\to0$, \eqref{eq:cert} forces $2wq\ge C>1$. Hence $2w(1-c)\ge1$ is necessary: if $2w(1-c)<1$, then \eqref{eq:cert} fails at $\mu_\delta$ for all small $\delta$.

\emph{Inequality \eqref{eq:D}.} Fix $x\in[\frac12,\frac1{\sqrt2}]$ with $\ent(x)\le w\,\ent(x^2)+1-w$, so that $p:=\ent(x)/(w\,\ent(x^2)+1-w)\in(0,1]$, and let $\mu_p=p\,\delta_x+(1-p)\,\delta_0$. Present the product law to every class. The i.i.d.\ term equals $p^2\,\ent(x^2)$ and every other term is at most $p^2$, so \eqref{eq:cert} at $\mu_p$ forces $p\,(w\,\ent(x^2)+1-w)\ge C\,\ent(x)>\ent(x)$, which fails at the chosen $p$ and at every slightly smaller $p$. If $c>1-x\,\ent(x)/(w\,\ent(x^2)+1-w)$ then $px>1-c$, so a slightly smaller $p$ still has $\E_{\mu_p}[X]=px\ge1-c$ and gives a law violating \eqref{eq:cert}. This proves \eqref{eq:D}. (The condition on $x$ holds at the binding point below, and restricting to it does not change the value of $c^{**}$.)

\emph{The fixed point.} The right side of \eqref{eq:D} is decreasing in $w$, because its denominator $1-w(1-\ent(x^2))$ is decreasing in $w$ and $\ent(x^2)\le1$. By \eqref{eq:H}, $w\ge w_0(c):=1/(2(1-c))$; substituting $w_0(c)$ into \eqref{eq:D} and minimising over $x\in X_c$ gives $c\le G(c)$. For $c\in[0.2,0.45]$ the constraint defining $X_c$ is inactive at the minimiser (the minimiser stays near $0.69$ while $X_c\supseteq[0.62,1/\sqrt2]$), so $G(c)$ coincides with the minimum over all of $[\frac12,\frac1{\sqrt2}]$; since $w_0$ is increasing in $c$, that minimum is decreasing in $c$, so $c\mapsto c-G(c)$ is strictly increasing on this range and $\{c:c\le G(c)\}=(0,c^{**}]$ with $c^{**}$ the unique root of $c=G(c)$ (for $c$ outside $[0.2,0.45]$ the inequality $c\le G(c)$ is checked directly). Its numerical value is computed in Appendix~\ref{app:const}.
\end{proof}

\begin{remark}\label{rem:literature}
(i) Cambie's exact value $0.3823455$ for Sawin's class lies below $c^{**}$ because the class of all couplings also lets the adversary reduce the mass of the diagonal pair $(x,x)$ from $p^2$ to $\max(0,2p-1)$, and Cambie's extremal law saturates the independent and the anti-correlated terms simultaneously. (ii) Liu's $0.382709$ lies below $c^{**}$ because the diagonal $x^2+x^2(1-x)^2$ of his kernel is not $\frac12$ at the binding point: his adversary is the two-point law of \eqref{eq:D} placed at the unique $x$ with $x^2+x^2(1-x)^2=1-x^2$ (namely $x=0.6908$, where the diagonal equals $0.5228$), so that $\ent$ of the kernel term equals $\ent(x^2)$ and the kernel is entropy-neutral there; with his weight $w=0.899947$ the resulting bound is exactly $0.382709$. Placing the two-point adversary against the optimal diagonal instead gives, for the same $w$, $0.382752$.
\end{remark}

\section{A protocol with the optimal diagonal}\label{sec:protocol}
\begin{definition}\label{def:ideal}
Let
\[f(x)=\begin{cases}\min(x,1-x),&0\le x\le\frac12,\\ \sqrt{\frac12-x^2},&\frac12\le x\le\frac1{\sqrt2},\\ 0,&\frac1{\sqrt2}\le x\le1,\end{cases}\]
and let $\Pi^{\id}$ be the conditionally-i.i.d.\ protocol with $r_x(u)=x+f(x)\,\mathrm{sgn}(u-\tfrac12)$, so that $\Pi^{\id}_{x,y}(0,0)=xy+f(x)f(y)$.
\end{definition}

\begin{lemma}\label{lem:ideal}
$\Pi^{\id}$ is a protocol. For all $x,y\in[0,1]$, $\ent(xy+f(x)f(y))\ge\ent(xy)$, and $x^2+f(x)^2=\frac12$ for $x\in[\frac12,\frac1{\sqrt2}]$.
\end{lemma}
\begin{proof}
We need $r_x\in[0,1]$, i.e.\ $f(x)\le\min(x,1-x)$. On $[0,\frac12]$ this is equality. On $[\frac12,\frac1{\sqrt2}]$, $\frac12-x^2\le(1-x)^2$ is equivalent to $2x^2-2x+\frac12\ge0$, i.e.\ $(2x-1)^2\ge0$. Next, $x^2+f(x)^2=2x^2\le\frac12$ on $[0,\frac12]$ and $=\frac12$ on $[\frac12,\frac1{\sqrt2}]$, so $x^2+f(x)^2\le\frac12$ wherever $f(x)>0$. If $f(x)f(y)>0$ then by Cauchy--Schwarz $xy+f(x)f(y)\le\sqrt{x^2+f(x)^2}\sqrt{y^2+f(y)^2}\le\frac12$; since $xy\le xy+f(x)f(y)\le\frac12$ and $\ent$ is increasing on $[0,\frac12]$, $\ent(xy+f(x)f(y))\ge\ent(xy)$. If $f(x)f(y)=0$ the two quantities coincide.
\end{proof}

\begin{proposition}[value against the two adversaries]\label{prop:twoadv}
Let $\Pi=w\,\Pi^{\iid}+(1-w)\,\Pi^{\id}$ denote the certificate \eqref{eq:cert} with the i.i.d.\ protocol at weight $w$ and $\Pi^{\id}$ (class $\cC_3$) at weight $1-w$.
\begin{enumerate}
\item[(a)] On the two-point law $\mu=p\,\delta_x+(1-p)\,\delta_0$, $x\in[\frac12,\frac1{\sqrt2}]$, the infimum over $\cC_3(\mu)$ of the $\Pi^{\id}$-term is attained at the product law and equals $p^2$; hence the ratio of the two sides of \eqref{eq:cert} is $p\,(w\,\ent(x^2)+1-w)/\ent(x)$, and \eqref{eq:D} holds with equality in the sense that the supremum of admissible $c$ against two-point laws is exactly
\[c_2(w):=\min_{x\in[1/2,1/\sqrt2]}\Big[1-\frac{x\,\ent(x)}{w\,\ent(x^2)+1-w}\Big].\]
\item[(b)] On the hiding laws $\mu_{q,y,\delta}$ with $q=1-c$, the ratio tends to $2w(1-c)$ as $\delta\to0$, from above.
\item[(c)] Let $w^{**}=1/(2(1-c^{**}))=0.810222\ldots$. Then $c_2(w^{**})=c^{**}$: at $w=w^{**}$ both adversaries are tight simultaneously, and for every $c<c^{**}$ there is $w$ with $2w(1-c)>1$ and $c<c_2(w)$.
\end{enumerate}
\end{proposition}
\begin{proof}
(a) The mixtures of products with marginal $\mu$ are the laws \eqref{eq:twoprod}; the $\Pi^{\id}$-term assigns weight only to the pair $(x,x)$ and there equals $\ent(\frac12)=1$ by Lemma~\ref{lem:ideal}. Its mass is $(1-q)P_0(x)^2+qP_1(x)^2\ge\big((1-q)P_0(x)+qP_1(x)\big)^2=p^2$ by convexity, with equality at the product law. The rest is the computation in the proof of Theorem~\ref{thm:refined}. (b) With the hiding coupling of Lemma~\ref{lem:classes} the $\Pi^{\id}$-term is $\delta^2(1-q)\,\ent(y^2+f(y)^2)\ge0$ and the i.i.d.\ term is $2q(1-q)\delta\ent(y)+(1-q)^2\delta^2\ent(y^2)$; dividing by $(1-q)\delta\ent(y)$ gives $2wq+O(\delta)$ with a non-negative $O(\delta)$ term. (c) By definition of $c^{**}$ as the root of $c=G(c)$ and $G(c)=c_2(w_0(c))$. For $c<c^{**}$ take $w$ strictly between $w_0(c)$ and $w_0(c^{**})$: then $2w(1-c)>1$ and $c_2(w)>c_2(w_0(c^{**}))=c^{**}>c$ by monotonicity of $c_2$ in $w$.
\end{proof}

Proposition~\ref{prop:twoadv} shows that the mixture $w\,\Pi^{\iid}+(1-w)\,\Pi^{\id}$ meets the necessary conditions of Theorem~\ref{thm:refined} with equality. What remains is to verify \eqref{eq:cert} against \emph{all} laws in $\cC_3$, i.e.\ against all $(q,P_0,P_1)$ in \eqref{eq:twoprod}. This is a computation, once the laws $P_0,P_1$ have been reduced to finitely many atoms.

\section{Cardinality reduction and the small-entropy regime}\label{sec:card}
Fix $w$, $c$ and $C>1$, and for $q\in[0,1]$ and probability laws $P_0,P_1$ on $[0,1]$ with $\mu=(1-q)P_0+qP_1$ define
\begin{equation}\label{eq:Phi}
\begin{aligned}
\Phi_C(q,P_0,P_1)&:=w\iint\ent(xy)\,d\mu\,d\mu+(1-w)\Big[(1-q)J(P_0)+qJ(P_1)\Big]-C\int\ent\,d\mu,\\
J(P)&:=\iint\ent\big(xy+f(x)f(y)\big)\,dP\,dP.
\end{aligned}
\end{equation}
By Liu's Theorem~9, the certificate \eqref{eq:cert} for $w\Pi^{\iid}+(1-w)\Pi^{\id}$ holds with constant $C$ if and only if $\Phi_C\ge0$ whenever $\E_\mu[X]\ge1-c$.

\begin{lemma}[extreme points of moment sets; \cite{Karr,Winkler}]\label{lem:extreme}
Let $g_1,\dots,g_r$ be continuous real functions on $[0,1]$ and $c_1,\dots,c_r\in\R$. The set $\cM=\{P\text{ probability law on }[0,1]:\int g_j\,dP=c_j,\ j=1,\dots,r\}$ is convex and weakly compact, and each of its extreme points is supported on at most $r+1$ points.
\end{lemma}

\begin{lemma}[concavity on slices]\label{lem:concave}
Let $\phi$ be a continuous symmetric function on $[0,1]^2$ and $g_1,\dots,g_r$ continuous on $[0,1]$. Suppose that $\iint\phi\,d\nu\,d\nu\le0$ for every finite signed measure $\nu$ on $[0,1]$ with $\int d\nu=0$ and $\int g_j\,d\nu=0$ for all $j$. Then $P\mapsto\iint\phi\,dP\,dP$ is concave on every set $\cM$ of Lemma~\ref{lem:extreme}.
\end{lemma}
\begin{proof}
For $P,Q\in\cM$ and $t\in[0,1]$, expanding the bilinear form gives $\iint\phi\,d(tP+(1-t)Q)^{\otimes2}-t\iint\phi\,dP^{\otimes2}-(1-t)\iint\phi\,dQ^{\otimes2}=-t(1-t)\iint\phi\,d(P-Q)\,d(P-Q)\ge0$, since $\nu=P-Q$ satisfies the hypotheses.
\end{proof}

The i.i.d.\ term satisfies the hypothesis of Lemma~\ref{lem:concave} with $\phi(x,y)=\ent(xy)$ and the single function $g_1(x)=x$: this is the concavity lemma of Alweiss, Huang and Sellke \cite[Lemma~5]{AHS} (see also \cite[proof of Theorem~9]{Liu}). For the kernel term we need the analogous statement, which for Liu's kernel he proves for small perturbations of the i.i.d.\ kernel and verifies numerically for his actual kernel \cite[Section~V-A]{Liu}. For $\Pi^{\id}$ the form is not negative semidefinite on $\{\int d\nu=\int x\,d\nu=0\}$: it has exactly two positive directions.

\begin{hypothesis}[inertia]\label{hyp:inertia}
There are continuous functions $g_2,g_3$ on $[0,1]$ such that $\iint\ent(xy+f(x)f(y))\,d\nu\,d\nu\le0$ for every finite signed measure $\nu$ with $\int d\nu=\int x\,d\nu=\int g_2\,d\nu=\int g_3\,d\nu=0$.
\end{hypothesis}
Numerically, discretising $[0,1]$ with step $10^{-3}$ and $5\cdot10^{-4}$ and projecting the matrix $[-\ent(x_ix_j+f(x_i)f(x_j))]_{ij}$ onto the orthogonal complement of the constant and linear vectors, the projected form has exactly two negative eigenvalues (stable across the two grids), separated from the remaining spectrum by six orders of magnitude; one of them is the same near-boundary direction that appears for Liu's kernel (where it is spanned, together with the constant and linear vectors, by Liu's moment function $x(1-x)$), and the other is concentrated on $x\in[0.696,0.705]$. Taking $g_2,g_3$ to be the interpolations of the two eigenvectors, the form projected onto the complement of $\{1,x,g_2,g_3\}$ has all eigenvalues above $-10^{-14}$, the level of rounding error, exactly as in Liu's check for his kernel. Hypothesis~\ref{hyp:inertia} is the continuous statement behind this computation; it has the same status as the positive-semidefiniteness hypothesis in \cite[Theorem~13]{Liu}.

\begin{proposition}[atom budget]\label{prop:atoms}
Assume Hypothesis~\ref{hyp:inertia}. Then for every $q\in[0,1]$, $\inf\Phi_C(q,P_0,P_1)$ over all pairs $(P_0,P_1)$ with $\E_\mu[X]\ge1-c$ equals the infimum over pairs in which $P_0$ and $P_1$ are each supported on at most four points.
\end{proposition}
\begin{proof}
Fix $q$, and fix the values $\int x\,dP_j$, $\int g_2\,dP_j$, $\int g_3\,dP_j$ for $j=0,1$; this fixes $\E_\mu[X]$. On the product $\cM_0\times\cM_1$ of the two corresponding slices, $J(P_0)$ and $J(P_1)$ are concave by Lemma~\ref{lem:concave} and Hypothesis~\ref{hyp:inertia}; the i.i.d.\ term is the composition of the linear map $(P_0,P_1)\mapsto\mu$, which maps $\cM_0\times\cM_1$ into a slice with fixed mean, with the concave functional of \cite[Lemma~5]{AHS}; and the last term of \eqref{eq:Phi} is linear. So $\Phi_C$ is concave and weakly continuous on the compact convex set $\cM_0\times\cM_1$, and by Bauer's minimum principle it attains its minimum at an extreme point, which is a pair of extreme points of $\cM_0$ and $\cM_1$; by Lemma~\ref{lem:extreme} with $r=3$ these have at most four atoms each. Every pair $(P_0,P_1)$ lies in some such product of slices.
\end{proof}

The remaining difficulty is the region of very small entropy, where the infimum of the ratio is approached along the hiding family rather than attained, so that no finite computation can exhaust it. The following bound handles it. Write $R(q,P_0,P_1)$ for the ratio of the two sides of \eqref{eq:cert} at a mixture of two products, and $\tilde L(u):=\log_2(1/u)+(1-u)/\ln2$, a decreasing function on $(0,1)$.

\begin{theorem}[small-entropy bound]\label{thm:small}
Let $w=0.810222$, $t_0=0.006$ and $\rho:=(1-t_0)/\ent(t_0)=18.7849$. For every law $\mu$ on $[0,1]$ with $\E_\mu[X]\ge1-c$ and $E:=\E_\mu\ent(X)>0$, and every $q,P_0,P_1$ with marginal $\mu$,
\begin{equation}\label{eq:small}
R(q,P_0,P_1)\;\ge\;2w\big[(1-c)-\rho E\big].
\end{equation}
In particular, for $c=0.38284$, $R\ge C=1.00005$ whenever $E\le\varepsilon_0:=7.6\cdot10^{-7}$.
\end{theorem}
The proof needs four elementary facts about $\ent$ and one exact identity.
\begin{enumerate}
\item[(F1)] $\ent(\theta u)\ge\theta\,\ent(u)$ for $\theta,u\in[0,1]$ (concavity, $\ent(0)=0$); equivalently $u\mapsto\ent(u)/u$ is decreasing on $(0,1]$.
\item[(F2)] $\ent(a)\le2\sqrt{a(1-a)}$ for $a\in[0,1]$.
\item[(F3)] $\ent(u)\ge u\,\tilde L(u)$ for $u\in(0,1)$: indeed $\ent(u)=u\log_2(1/u)+(1-u)\log_2(1/(1-u))$ and $\ln(1/(1-u))\ge u$.
\item[(F4)] For $x=1-t$, $y=1-t'$ with $t,t'\in[0,1]$, let $\Lambda(t,t'):=\ent(t)+\ent(t')-\ent(xy)$. Then $\Lambda\ge0$ and, with $\tau=1-xy=t+t'-tt'$, $A=t(1-t')$, $B=t'(1-t)$, $\gamma=tt'/\tau$, $\alpha=A/(A+B)$,
\[\Lambda(t,t')=\tau\,\ent(\gamma)+(A+B)\,\ent(\alpha).\]
Proof: for independent $X\sim\mathrm{Bern}(x)$, $Y\sim\mathrm{Bern}(y)$, $\ent(t)+\ent(t')=H(X,Y)=H(XY)+H(X,Y\mid XY)$, and given $XY=0$ (probability $\tau$) the pair takes the values $(0,0),(1,0),(0,1)$ with probabilities proportional to $tt',A,B$.
\item[(F5)] For $t,t'\in(0,\frac12]$ and $m=\min(t,t')$: $\Lambda(t,t')\le2\sqrt{tt'}+tt'\big(\tilde L(m)+1+m/\ln2\big)$. Proof: by (F2), $(A+B)\ent(\alpha)\le2\sqrt{AB}\le2\sqrt{tt'}$; and $\tau\ent(\gamma)=tt'\log_2(1/\gamma)+\tau(1-\gamma)\log_2(1/(1-\gamma))\le tt'\log_2\big(\tfrac1t+\tfrac1{t'}\big)+tt'/\ln2\le tt'\big(\log_2(2/m)+1/\ln2\big)$, using $\tau\le t+t'$ and $(1-\gamma)\ln(1/(1-\gamma))\le\gamma$.
\end{enumerate}
(F2) is checked as follows: for $a=\frac12-u$ with $u\le0.4612$, $\ent''\le-4/\ln2$ gives $\ent(a)\le1-2u^2/\ln2\le\sqrt{1-4u^2}$; for $a<0.0389$, $\ent(a)\le\sqrt a\cdot\sqrt a\,(\log_2(1/a)+1/\ln2)\le1.2077\sqrt a<1.9608\sqrt a\le2\sqrt{a(1-a)}$, since $\sqrt a(\log_2(1/a)+1/\ln2)$ is increasing on $(0,0.3679)$.

\begin{lemma}[pointwise inequality]\label{lem:pointwise}
For all $t,t'\in[0,t_0]$, $w\,\Lambda(t,t')\le(1-w)\sqrt{\ent(t)\ent(t')}$.
\end{lemma}
\begin{proof}
If $tt'=0$ then $\Lambda=0$. Otherwise let $m=\min(t,t')\le M=\max(t,t')\le t_0$. By (F3), $\sqrt{\ent(t)\ent(t')}\ge\sqrt{mM\,\tilde L(m)\tilde L(M)}$, so by (F5) it suffices that
\[\frac{2}{\sqrt{\tilde L(m)\tilde L(M)}}+\frac{\sqrt{mM}\,(\tilde L(m)+1+m/\ln2)}{\sqrt{\tilde L(m)\tilde L(M)}}\;\le\;r:=\frac{1-w}{w}=0.234230.\]
Since $\tilde L$ is decreasing, the first term is at most $2/\tilde L(t_0)$. The second term equals $\big[\sqrt{m\tilde L(m)}+\sqrt m(1+m/\ln2)/\sqrt{\tilde L(m)}\big]\cdot\sqrt{M/\tilde L(M)}$; each factor is increasing ($\frac{d}{dm}[m\tilde L(m)]=\log_2(1/m)-2m/\ln2>0$ on $(0,t_0]$), so it is at most its value at $m=M=t_0$. Altogether the left side is at most $2/\tilde L(t_0)+t_0(\tilde L(t_0)+1+t_0/\ln2)/\tilde L(t_0)=0.226890+0.006687=0.233577<r$.
\end{proof}

\begin{proof}[Proof of Theorem~\ref{thm:small}]
Let $N=(1-t_0,1]$, $B=[0,1-t_0]$, and write $t_x=1-x$ for $x\in N$, $\sigma=\mu(N)$, $S=\int_Nx\,d\mu\le\sigma$, $E_N=\int_N\ent\,d\mu$, $E_B=\int_B\ent\,d\mu$, $E=E_N+E_B$. Note $\ent(x)=\ent(t_x)$.

\emph{Kernel term.} For any mixture $\sum_j\pi_jP_j\otimes P_j$ of products with marginal $\mu$ (in particular for \eqref{eq:twoprod}), using $\ent(xy+f(x)f(y))\ge\ent(xy)\ge0$ and keeping only $N\times N$: for $x,y\in N$, $1-xy\in[\max(t_x,t_y),2t_0]\subset[0,\frac12]$, so $\ent(xy)\ge\ent(\max(t_x,t_y))\ge\sqrt{\ent(t_x)\ent(t_y)}$, whence $\iint\ent(xy+ff)\,dP_j\,dP_j\ge(\int_N\sqrt\ent\,dP_j)^2$, and by Jensen $\sum_j\pi_j(\int_N\sqrt\ent\,dP_j)^2\ge(\int_N\sqrt\ent\,d\mu)^2=\iint_{N\times N}\sqrt{\ent(t_x)\ent(t_y)}\,d\mu\,d\mu$.

\emph{I.i.d.\ term.} Drop $B\times B$. For $x\in N$, $y\in B$, (F1) gives $\ent(xy)\ge x\,\ent(y)$, so the two cross terms contribute at least $2SE_B$. On $N\times N$, $\ent(xy)=\ent(t_x)+\ent(t_y)-\Lambda(t_x,t_y)$ by definition, contributing $2\sigma E_N-\iint_{N\times N}\Lambda$.

\emph{Combination.} Therefore
\begin{align*}
w\iint\ent(xy)\,d\mu\,d\mu+(1-w)\sum_j\pi_jJ(P_j)\;&\ge\;2wSE_B+2w\sigma E_N+\iint_{N\times N}D(t_x,t_y)\,d\mu\,d\mu\\
&\ge\;2wSE_B+2w\sigma E_N\;\ge\;2wS\,E,
\end{align*}
where $D(t,t'):=(1-w)\sqrt{\ent(t)\ent(t')}-w\Lambda(t,t')$ is non-negative by Lemma~\ref{lem:pointwise} (all $t_x,t_y<t_0$), and $\sigma\ge S$.

\emph{Mean constraint.} By (F1), $y/\ent(y)$ is increasing on $(0,1]$, so $y\le\rho\,\ent(y)$ on $B$ with $\rho=(1-t_0)/\ent(1-t_0)=(1-t_0)/\ent(t_0)$. Hence $1-c\le\int x\,d\mu=S+\int_By\,d\mu\le S+\rho E_B\le S+\rho E$.

Dividing by $E$ gives $R\ge2wS\ge2w[(1-c)-\rho E]$, which is at least $C$ iff $E\le[(1-c)-C/(2w)]/\rho=1.4329\cdot10^{-5}/18.7849=7.628\cdot10^{-7}$.
\end{proof}

Two features make \eqref{eq:small} sharp enough. A near-$1$ atom $x=1-t$ contributes $x\,\ent(y)$ to the cross terms, and the same factor $x$ enters the mean constraint through $S$, so the deficiency of such an atom cancels its cross-term loss exactly; and inside $N$ the subadditivity loss $\Lambda\approx2\sqrt{tt'}$ is paid for exactly by the kernel term, with no factor $1/\log_2(1/t)$ lost. The one genuine loss is the mean-deficiency step, which charges mass just below $1-t_0$ at rate $\rho$; this produces the linear term and hence $\varepsilon_0$. Numerically, the infimum of $R$ over laws with $E\le10^{-3}$ appears to be exactly $2w(1-c)$, so $\varepsilon_0$ is conservative.

\section{Numerical certification and the main theorem}\label{sec:numerics}
By Proposition~\ref{prop:atoms}, certifying $\Phi_C\ge0$ reduces, under Hypothesis~\ref{hyp:inertia}, to a minimisation over $q\in[0,1]$ and two laws with at most four atoms each: the variables are $q$, the atoms $x_{j,1},\dots,x_{j,4}\in[0,1]$ and weights of $P_j$ ($j=0,1$), 17 real parameters, with the constraint $\E_\mu[X]\ge1-c$. We minimised the ratio
\begin{equation}\label{eq:ratio}
R(q,P_0,P_1):=\frac{w\iint\ent(xy)\,d\mu\,d\mu+(1-w)\big[(1-q)J(P_0)+qJ(P_1)\big]}{\int\ent\,d\mu}
\end{equation}
(so $\Phi_C\ge0$ iff $R\ge C$ whenever $\int\ent\,d\mu>0$) with sequential least-squares programming from $300$--$500$ random starting points per run ($5{,}000$ for the decisive cell, four atoms at $c=0.38284$), together with structured starting points: two-point laws $p\,\delta_x+(1-p)\delta_0$ for $x\in[0.6,0.75]$; the hiding laws \eqref{eq:mudelta} with $\delta\in\{10^{-2},10^{-3},10^{-4}\}$; the extremal laws of Cambie and Liu; and combinations of the two-point and hiding laws. Because the infimum along the hiding family is approached only as $\delta\to0$, where the denominator of $R$ vanishes, that family was evaluated in closed form (Proposition~\ref{prop:twoadv}(b)), and the optimiser was run with the additional constraint $\int\ent\,d\mu\ge\eta_0$ for $\eta_0=10^{-3},10^{-4},10^{-5},10^{-6},10^{-7}$, always with the same minimum and minimiser. The atom budget was also increased to five and six per law as a check. The computation was implemented independently twice (the second time from the statement alone, without access to the first implementation), and the evaluator reproduces Liu's optimum for his kernel to nine digits: $R=1.000000000$ at his $(c',w)$, attained at his minimiser $q=0$, $P_0=0.8936\,\delta_{0.6908}+0.1064\,\delta_0$.

\begin{table}[h]\centering\footnotesize
\begin{tabular}{lcrr}\toprule
adversary family & atoms per law & $c=0.38284$ & $c=0.38288$\\\midrule
all laws, $\int\ent\,d\mu\ge10^{-3}$, 300--500 starts & 4 & 1.0000733 & 1.0000085\\
 & 5 & 1.0000733 & 1.0000085\\
 & 6 & 1.0000733 & 1.0000085\\
hiding laws \eqref{eq:mudelta}, $\delta\le10^{-2}$, four values of $y$, closed form & --- & $\ge1.000073$ & ---\\
two-point/hiding combinations, seeded & 4 & 1.0000733 & ---\\\bottomrule
\end{tabular}
\caption{Minimum of $R$ in \eqref{eq:ratio} at $w=0.810222$. Every minimiser found is the two-point law $0.893\,\delta_{0.6909}+0.107\,\delta_0$ (with $q\in\{0,1\}$), whose value is the closed-form expression of Proposition~\ref{prop:twoadv}(a); the hiding family decreases to its limit $2w(1-c)=1.000073$.}\label{tab:cert}
\end{table}

\begin{hypothesis}[global minimum]\label{hyp:global}
For $w=0.810222$ and $c=0.38284$, the minimum of $R$ over $q\in[0,1]$ and laws $P_0,P_1$ with at most four atoms each, $\E_\mu[X]\ge1-c$ and $\int\ent\,d\mu\ge10^{-7}$, is $1.0000733$, attained at the two-point law of Table~\ref{tab:cert}; in particular $R\ge C=1.00005$ on this set.
\end{hypothesis}
Hypothesis~\ref{hyp:global} is supported by Table~\ref{tab:cert} and by the same minima with the floor lowered from $10^{-3}$ to $10^{-7}$ (in every case the same minimiser and the same value). The complementary region $\int\ent\,d\mu<10^{-7}$ is covered unconditionally by Theorem~\ref{thm:small}, whose threshold $\varepsilon_0=7.6\cdot10^{-7}$ exceeds the floor, so the two regions overlap and no window remains. This is the point at which the present certification is more complete than Liu's, whose analogue of Theorem~\ref{thm:small} \cite[Lemma~7]{Liu} is a statement about limits and whose search used no entropy floor. The hypothesis otherwise plays the role of the ``global minimiser structure'' hypothesis in \cite[Theorem~13]{Liu}. The analogous hypothesis at other $(w,c)$, with $C$ replaced by any constant strictly between $1$ and the two-point value of Proposition~\ref{prop:twoadv}(a) and with the threshold of Theorem~\ref{thm:small} recomputed as $\varepsilon_0(w,c,C)=[2w(1-c)-C]/(2w\rho)$, is denoted Hypothesis~\ref{hyp:global}$(w,c)$; we verified it, with the same minimiser, at $(0.810222,0.38288)$ as well, where the two-point value is $1.0000085$, $2w(1-c)=1.0000072$, and the best threshold is $2.8\cdot10^{-7}$, still above the floor. Since $2w(1-c)-C$ vanishes as $c\to c^{**}$, the threshold falls below the floor $10^{-7}$ once $c>c^{**}-1.9\cdot10^{-6}=0.3828834$; for such $c$ a shrinking small-entropy window is again hypothesised rather than proved.

\begin{theorem}\label{thm:main}
Assume Hypotheses~\ref{hyp:inertia} and~\ref{hyp:global}. Then the certificate \eqref{eq:cert} holds for $w\Pi^{\iid}+(1-w)\Pi^{\id}$ with $w=0.810222$, $c=0.38284$ and $C=1.00005$; consequently every finite union-closed family $\F\ne\{\emptyset\}$ has an element contained in at least $0.38284\,|\F|$ of its members. For any $c<c^{**}=0.382885$, the same conclusion holds under Hypothesis~\ref{hyp:inertia} and Hypothesis~\ref{hyp:global}$(w,c)$ with $w$ chosen as in Proposition~\ref{prop:twoadv}(c); for $c\le c^{**}-1.9\cdot10^{-6}$ the small-entropy regime is covered by Theorem~\ref{thm:small}, as at $c=0.38284$.
\end{theorem}
\begin{proof}
Let $\mu$ have $\E_\mu[X]\ge1-c$ and let $P\in\cC_3(\mu)$. By Liu's Theorem~9 we may assume $P$ is of the form \eqref{eq:twoprod}. If $\int\ent\,d\mu=0$ both sides of \eqref{eq:cert} vanish. Otherwise, by Proposition~\ref{prop:atoms} (with $C=1.00005$), $\Phi_C(q,P_0,P_1)\ge\Phi_C(q,P_0',P_1')$ for some pair $(P_0',P_1')$ with at most four atoms each on the same slice (so with the same $\E_\mu[X]$). If $\int\ent\,d\mu'\ge10^{-7}$, Hypothesis~\ref{hyp:global} gives $R\ge C$ there; if $0<\int\ent\,d\mu'<10^{-7}<\varepsilon_0$, Theorem~\ref{thm:small} gives $R\ge C$ (note $\E_{\mu'}[X]=\E_\mu[X]\ge1-c$); in both cases $\Phi_C(q,P_0',P_1')\ge0$, and if $\int\ent\,d\mu'=0$ then $\Phi_C(q,P_0',P_1')\ge0$ trivially. Proposition~\ref{prop:cert} concludes. The final sentence follows in the same way from Proposition~\ref{prop:twoadv}(c) and Hypothesis~\ref{hyp:global}$(w,c)$.
\end{proof}

\begin{remark}
Two features distinguish this certification from Liu's. First, the cardinality reduction needs two auxiliary functionals rather than one (Hypothesis~\ref{hyp:inertia}), because the kernel of $\Pi^{\id}$ is far from the i.i.d.\ kernel; the price is four atoms per law instead of three. Second, the hiding family, which is exactly what limits the weight $1-w$ of the kernel protocol, is handled analytically (Proposition~\ref{prop:twoadv}(b) and Theorem~\ref{thm:small}); numerical minimisation alone cannot resolve it, and a first, uncorrected computation of ours reported values up to $0.38295$ for larger kernel weights that are refuted by \eqref{eq:H}.
\end{remark}

\section{Discussion}\label{sec:discussion}
Theorem~\ref{thm:refined} says that the single-letter method, in every form used since 2022, is exhausted at $c^{**}=0.382885$, and Theorem~\ref{thm:main} says that the last $1.8\cdot10^{-4}$ above Liu's constant is available. The two adversaries explain the shape of the record. The diagonal law of Theorem~\ref{thm:ceil} is what Gilmer's argument was always fighting, and the successive improvements are successive recoveries of the diagonal deficit $1-\ent(x^2)$. The hiding law is what prevents the recovery from being complete: because the certificate must hold for every law $\mu$, and in particular for laws in which almost all entropy sits on an atom of vanishing mass, any protocol whose class allows the adversary to keep that atom away from the ``surely absent'' atom receives none of the entropy, and only the i.i.d.\ term, whose class is a singleton, is guaranteed the cross pairs. Hence $w\ge1/(2(1-c))$.

Going beyond $c^{**}$ therefore requires one of two things. Either a \emph{tensorising} constraint on the joint law of the two prefix-conditional probabilities that forbids hiding a tiny atom---maximal correlation does not, as Lemma~\ref{lem:classes} shows, because a tiny atom can be moved at vanishing cost in correlation---or an argument that is not single-letter, for instance one that uses the exact closure hypothesis through $H(A\cup B)\le\log_2|\F|-D\big(\mathrm{law}(A\cup B)\,\|\,\mathrm{Unif}(\F)\big)$ rather than through \eqref{eq:union}. Some evidence that the loss is in the single-letter relaxation rather than in the coupling idea comes from an exact dynamic-programming oracle for prefix-dependent (non-memoryless) couplings on small families: on every family on at most four elements with maximal frequency below $\frac12$, and on all unions of two Hamming layers on $24$ elements, some sequential coupling makes the lower bound in \eqref{eq:chain} exceed $H(A)$. We do not know how to turn this into a proof, because the adversary in the single-letter relaxation chooses the joint law of the two prefixes, which the dynamic programme controls.

\section*{Acknowledgements}
The literature survey, the transcription of the cited proofs, the numerical computations, including the independent re-implementation of the certificate, and the proof of Theorem~\ref{thm:small} were carried out with Claude (Anthropic) language-model agents under the author's direction; the mathematical claims were checked by the author and, separately, over several rounds by an agent instructed to referee the manuscript. Code, logs, the referee reports and the literature transcripts are available at \url{https://github.com/moffatstudio/union-closed-constant}. The Lean formalisation of Appendices~\ref{app:lean} and~\ref{app:src} was written by a Claude (Anthropic) agent from a written specification and checked against the paper by the author.

\appendix
\section{Constants}\label{app:const}
With $\ent$ in bits, $\ent(1/\sqrt2)=0.872429340\ldots$ and $c_{\mathrm{ceil}}=1-\ent(1/\sqrt2)/\sqrt2=0.383099298\ldots$. The fixed point of \eqref{eq:G} is
\[c^{**}=0.382885260\ldots,\qquad w^{**}=\frac1{2(1-c^{**})}=0.810222\ldots,\qquad x^{**}=0.690908\ldots,\]
where $x^{**}$ is the minimiser in \eqref{eq:G}, with $\ent(x^{**})=0.892123$, $\ent(x^{**2})=0.998520$ and $x^{**2}+f(x^{**})^2=\frac12$. For comparison, $c_2(w)$ of Proposition~\ref{prop:twoadv}(a) equals $0.382752$ at Liu's weight $w=0.899947$ (minimiser $x=0.6807$), $0.382841$ at $w=0.85$, and $0.382885$ at $w=w^{**}$; the constraint \eqref{eq:H} at $c=0.38284$ reads $w\ge0.81016$, and at $w=0.810222$ the hiding family gives $2w(1-c)=1.000073$. The maximal correlation of the hiding coupling of Lemma~\ref{lem:classes} at $q=1-0.38284$, $y=\frac12$ is $0.081$, $0.025$, $0.0079$, $0.00079$ for $\delta=10^{-2},10^{-3},10^{-4},10^{-6}$. In Theorem~\ref{thm:small}: $\tilde L(t_0)=8.814861$, $\ent(t_0)=0.0529151$, $\rho=18.784815$, $r=(1-w)/w=0.2342296$, $(1-c)-C/(2w)=1.43288\cdot10^{-5}$, $\varepsilon_0=7.62787\cdot10^{-7}$; the ratio $w\Lambda/((1-w)\sqrt{\ent\ent'})$ on $[0,t_0]^2$ is largest at the corner, where it equals $0.9909$, and Lemma~\ref{lem:pointwise} holds up to $t_0=0.00631$.


\section{Formal verification of the ceiling theorems}\label{app:lean}
\sloppy
Theorem~\ref{thm:ceil}, Theorem~\ref{thm:refined} and the three formalised cases of
Lemma~\ref{lem:classes} have been checked in Lean~4 with Mathlib. The development is a
single Lean library, \texttt{UnionClosedCeiling}, of about $1{,}040$ lines; it contains no
\texttt{sorry}, and every main theorem depends only on Lean's three standard axioms. This
appendix describes the model, quotes the formal statements, says what is \emph{not}
formalised, and records the certificate of the check. Appendix~\ref{app:src} reproduces the
complete source.

\subsection{The finite model}\label{app:lean:model}
A law $\mu$ on $[0,1]$ is represented by finitely many atoms $a:\mathrm{Fin}\,n\to\R$ in
$[0,1]$ with weights $m:\mathrm{Fin}\,n\to\R$, $m\ge0$, $\sum_im_i=1$ (\texttt{IsLaw}), with
$\E_\mu[X]=\sum_im_ia_i$ (\texttt{mean}) and $\E_\mu[\ent(X)]=\sum_im_i\ent(a_i)$
(\texttt{entropy}). A coupling of $\mu$ with itself is a non-negative matrix whose row sums
and column sums are $m$ (\texttt{IsCoupling}); the product law is
$P_{ij}=m_im_j$ (\texttt{productCoupling}). A protocol is represented by the single number
$\Pi_{x,y}(0,0)$, i.e.\ by a function $\Pi:\R\to\R\to\R$ subject to the Fr\'echet bounds
$\max(0,x+y-1)\le\Pi(x,y)\le\min(x,y)$ on $[0,1]^2$ (\texttt{IsProtocol}); Section~\ref{sec:ceiling}
uses nothing else about protocols, and the facts it does use---$\Pi_{x,0}(0,0)=0$,
$\Pi_{0,y}(0,0)=0$, $\Pi_{1,1}(0,0)=1$, $\Pi_{1,y}(0,0)=y$ (\texttt{protocol\_right\_zero}, \texttt{protocol\_left\_zero}, \texttt{protocol\_one\_one}, \texttt{protocol\_one\_left})---are derived from those bounds, while $\ent(\Pi)\le1$ holds for every real argument (\texttt{h\_le\_one}, from Mathlib's \texttt{Real.binEntropy\_le\_log\_two})
in the Lean file. A class is a predicate on (arity, atoms, weights, coupling), and the two
properties of classes used in Section~\ref{sec:ceiling} are \texttt{ContainsProduct} and
\texttt{AdmitsHiding}, the latter being Definition~\ref{def:hiding}: for all $q,y\in(0,1)$
there is a threshold $\delta_0>0$ such that for every $\delta\in(0,\delta_0)$ the class
contains a coupling $P$ of $\mu_{q,y,\delta}$ with itself with $P(y,1)=P(1,y)=0$ and
$P(y,y)\le\delta^2$. The two adversaries are the finitely supported laws
$p\,\delta_x+(1-p)\delta_0$ (\texttt{twoAtoms}, \texttt{twoWeights}) and $\mu_{q,y,\delta}$
(\texttt{hidingAtoms}, \texttt{hidingWeights}), indexed $0\mapsto1$, $1\mapsto0$,
$2\mapsto y$. Entropy is Mathlib's \texttt{Real.binEntropy} divided by $\log2$, so that
$\ent(\tfrac12)=1$.

The hypothesis of every theorem is \texttt{Certifies}: for every finitely supported law of
mean at least $1-c$ and every family of couplings $(P^{(k)})_k$ with $P^{(k)}$ admissible for
the $k$-th class,
$C\,\E_\mu[\ent(X)]\le\sum_kw_k\,\E_{P^{(k)}}\big[\ent(\Pi^{(k)}_{X,Y}(0,0))\big]$.
This is implied by \eqref{eq:cert}: the infimum over a class is at most the value at any
admissible family, so \eqref{eq:cert} yields \texttt{Certifies} for every such family, and
every Lean theorem therefore applies verbatim to the paper's certificate. Restricting to
finitely supported laws also only weakens the hypothesis, hence strengthens the theorems
proved from it; every adversary of Section~\ref{sec:ceiling} is of that form (two atoms in
Theorem~\ref{thm:ceil} and in \eqref{eq:D}, three in Definition~\ref{def:hiding}). Where the
formal hypotheses differ from the paper's they are weaker: the Lean statements assume
$C\ge1$ where the paper has $C>1$, and they give the i.i.d.\ protocol a class satisfying
only \texttt{ContainsProduct}, whereas the paper gives it the singleton class
$\{\mu\otimes\mu\}$, which satisfies that and more. In Theorem~\ref{thm:refined} the
protocols are indexed by $\mathrm{Fin}(K+1)$ and index $0$ is the i.i.d.\ one, characterised
by the hypothesis $\Pi^{(0)}_{x,y}(0,0)=xy$. The limit $\delta\to0$ in the proof of
\eqref{eq:H} is carried out by hand: $\delta$ is chosen below the hiding threshold of each of
the finitely many classes and below $c(1-2qw_0)$, which yields the contradiction directly.

\subsection{The statements}\label{app:lean:stmts}
The following are copied from the sources; only the statements are shown, the proofs are in
Appendix~\ref{app:src}. Definitions, from \texttt{UnionClosedCeiling/\allowbreak Framework.lean}, lines
61--63, 122--128, 131--133 and 205--211:

\begin{Verbatim}[fontsize=\footnotesize,formatcom=\leanmono,samepage=false]
def IsProtocol (Prot : ℝ → ℝ → ℝ) : Prop :=
  ∀ x y, x ∈ Set.Icc (0:ℝ) 1 → y ∈ Set.Icc (0:ℝ) 1 →
    max 0 (x + y - 1) ≤ Prot x y ∧ Prot x y ≤ min x y
\end{Verbatim}
\begin{Verbatim}[fontsize=\footnotesize,formatcom=\leanmono,samepage=false]
def Certifies (Prot : Fin K → ℝ → ℝ → ℝ) (w : Fin K → ℝ)
    (C : Fin K → ∀ n : ℕ, (Fin n → ℝ) → (Fin n → ℝ) → (Fin n → Fin n → ℝ) → Prop)
    (Cc c : ℝ) : Prop :=
  ∀ (n : ℕ) (a m : Fin n → ℝ), IsLaw a m → 1 - c ≤ mean a m →
    ∀ P : Fin K → (Fin n → Fin n → ℝ),
      (∀ k, IsCoupling m (P k) ∧ C k n a m (P k)) →
      Cc * entropy a m ≤ ∑ k, w k * expect a (P k) (fun x y => h (Prot k x y))
\end{Verbatim}
\begin{Verbatim}[fontsize=\footnotesize,formatcom=\leanmono,samepage=false]
def ContainsProduct (C : ∀ n : ℕ, (Fin n → ℝ) → (Fin n → ℝ) → (Fin n → Fin n → ℝ) → Prop) :
    Prop :=
  ∀ (n : ℕ) (a m : Fin n → ℝ), IsLaw a m → C n a m (productCoupling m)
\end{Verbatim}
\begin{Verbatim}[fontsize=\footnotesize,formatcom=\leanmono,samepage=false]
def AdmitsHiding (C : ∀ n : ℕ, (Fin n → ℝ) → (Fin n → ℝ) → (Fin n → Fin n → ℝ) → Prop) :
    Prop :=
  ∀ q y : ℝ, q ∈ Set.Ioo (0:ℝ) 1 → y ∈ Set.Ioo (0:ℝ) 1 →
    ∃ δ₀ > (0:ℝ), ∀ δ : ℝ, 0 < δ → δ < δ₀ →
      ∃ P : Fin 3 → Fin 3 → ℝ, IsCoupling (hidingWeights q δ) P ∧
        C 3 (hidingAtoms y) (hidingWeights q δ) P ∧
        P 2 0 = 0 ∧ P 0 2 = 0 ∧ P 2 2 ≤ δ^2
\end{Verbatim}

Theorem~\ref{thm:ceil} is \texttt{product\_ceiling}, from
\texttt{UnionClosedCeiling/\allowbreak Ceiling.lean}, lines 36--42 (the conclusion is the
contrapositive of the statement in the text):

\begin{Verbatim}[fontsize=\footnotesize,formatcom=\leanmono,samepage=false]
theorem product_ceiling {K : ℕ} {Prot : Fin K → ℝ → ℝ → ℝ} {w : Fin K → ℝ}
    {C : Fin K → ∀ n : ℕ, (Fin n → ℝ) → (Fin n → ℝ) → (Fin n → Fin n → ℝ) → Prop}
    {Cc c : ℝ}
    (hProt : ∀ k, IsProtocol (Prot k)) (hw : ∀ k, 0 ≤ w k) (hw1 : ∑ k, w k = 1)
    (hC : ∀ k, ContainsProduct (C k)) (hCc : 1 ≤ Cc)
    (hcert : Certifies Prot w C Cc c) :
    c ≤ 1 - h (1 / Real.sqrt 2) / Real.sqrt 2
\end{Verbatim}

The two inequalities of Theorem~\ref{thm:refined}, its fixed-point form and its numerical
consequence are \texttt{hiding\_bound} (\eqref{eq:H}), \texttt{diagonal\_bound}
(\eqref{eq:D}), \texttt{fixed\_point\_form} (the substitution of \eqref{eq:H} into
\eqref{eq:D}, i.e.\ \eqref{eq:G}) and \texttt{refined\_ceiling\_numeric}, from
\texttt{UnionClosedCeiling/\allowbreak Refined.lean}, lines 158--164, 287--292, 373--379 and 420--425. Two statements are wider than in the paper: \texttt{diagonal\_bound} is stated for $x\in(0,1]$ and \texttt{fixed\_point\_form} for $x\in(0,1)$, whereas Theorem~\ref{thm:refined} restricts $x$ to $[\tfrac12,\tfrac1{\sqrt2}]$; the proofs never use that restriction, so the Lean statements imply the paper's.
They share the section variables

\begin{Verbatim}[fontsize=\footnotesize,formatcom=\leanmono,samepage=false]
variable {K : ℕ} {Prot : Fin (K+1) → ℝ → ℝ → ℝ} {w : Fin (K+1) → ℝ}
  {C : Fin (K+1) → ∀ n : ℕ, (Fin n → ℝ) → (Fin n → ℝ) → (Fin n → Fin n → ℝ) → Prop}
  {Cc c : ℝ}
\end{Verbatim}

\begin{Verbatim}[fontsize=\footnotesize,formatcom=\leanmono,samepage=false]
theorem hiding_bound
    (hProt : ∀ k, IsProtocol (Prot k)) (hiid : ∀ x y, Prot 0 x y = x * y)
    (hw : ∀ k, 0 ≤ w k) (hw1 : ∑ k, w k = 1)
    (hC0 : ContainsProduct (C 0)) (hChide : ∀ k, k ≠ 0 → AdmitsHiding (C k))
    (hCc : 1 ≤ Cc) (hc0 : 0 < c) (hc1 : c < 1)
    (hcert : Certifies Prot w C Cc c) :
    1 ≤ 2 * w 0 * (1 - c)
\end{Verbatim}
\begin{Verbatim}[fontsize=\footnotesize,formatcom=\leanmono,samepage=false]
theorem diagonal_bound
    (hProt : ∀ k, IsProtocol (Prot k)) (hiid : ∀ x y, Prot 0 x y = x * y)
    (hw : ∀ k, 0 ≤ w k) (hw1 : ∑ k, w k = 1) (hC : ∀ k, ContainsProduct (C k))
    (hCc : 1 ≤ Cc) (hc1 : c < 1) (hcert : Certifies Prot w C Cc c)
    {x : ℝ} (hx : x ∈ Set.Ioc (0:ℝ) 1) (hcond : h x ≤ w 0 * h (x^2) + 1 - w 0) :
    c ≤ 1 - x * h x / (w 0 * h (x^2) + 1 - w 0)
\end{Verbatim}
\begin{Verbatim}[fontsize=\footnotesize,formatcom=\leanmono,samepage=false]
theorem fixed_point_form
    (hProt : ∀ k, IsProtocol (Prot k)) (hiid : ∀ x y, Prot 0 x y = x * y)
    (hw : ∀ k, 0 ≤ w k) (hw1 : ∑ k, w k = 1) (hC : ∀ k, ContainsProduct (C k))
    (hChide : ∀ k, k ≠ 0 → AdmitsHiding (C k))
    (hCc : 1 ≤ Cc) (hc0 : 0 < c) (hchalf : c ≤ 1/2) (hcert : Certifies Prot w C Cc c)
    {x : ℝ} (hx : x ∈ Set.Ioo (0:ℝ) 1) (hmono : h x ≤ h (x^2)) :
    c ≤ 1 - x * h x / (1 - (1 - h (x^2)) / (2 * (1 - c)))
\end{Verbatim}
\begin{Verbatim}[fontsize=\footnotesize,formatcom=\leanmono,samepage=false]
theorem refined_ceiling_numeric
    (hProt : ∀ k, IsProtocol (Prot k)) (hiid : ∀ x y, Prot 0 x y = x * y)
    (hw : ∀ k, 0 ≤ w k) (hw1 : ∑ k, w k = 1) (hC : ∀ k, ContainsProduct (C k))
    (hChide : ∀ k, k ≠ 0 → AdmitsHiding (C k))
    (hCc : 1 ≤ Cc) (hc0 : 0 < c) (hchalf : c ≤ 1/2) (hcert : Certifies Prot w C Cc c) :
    c ≤ 3829/10000
\end{Verbatim}

The five formalised statements of Lemma~\ref{lem:classes} (three classes: products for all three, hiding for two) are the following, from
\texttt{UnionClosedCeiling/\allowbreak Classes.lean}, lines 54, 51, 57, 131 and 141 in the order
shown, where \texttt{iidClass} is the singleton class
$\{\mu\otimes\mu\}$, \texttt{allCouplings} is Sawin's class and \texttt{mixtureOfProducts}
is the class $\cC_3$ of \eqref{eq:twoprod}:

\begin{Verbatim}[fontsize=\footnotesize,formatcom=\leanmono,samepage=false]
lemma containsProduct_iidClass : ContainsProduct iidClass
lemma containsProduct_allCouplings : ContainsProduct allCouplings
lemma containsProduct_mixtureOfProducts : ContainsProduct mixtureOfProducts
lemma admitsHiding_allCouplings : AdmitsHiding allCouplings
lemma admitsHiding_mixtureOfProducts : AdmitsHiding mixtureOfProducts
\end{Verbatim}

\subsection{What is not formalised}\label{app:lean:gaps}
Proposition~\ref{prop:cert} itself is not formalised: Gilmer's argument, the chain rule and
the reduction to union-closed families are outside the scope of the development, and the
certificate is a hypothesis there. The maximal-correlation classes of Lemma~\ref{lem:classes}
are not formalised either: their hiding witness rests on the continuity of singular values in
the matrix entries, a genuinely different argument. Sections~\ref{sec:protocol}--\ref{sec:numerics}
are outside the scope entirely: the protocol $\Pi^{\id}$, the cardinality reduction, the
small-entropy bound, the numerics and everything conditional on
Hypotheses~\ref{hyp:inertia} and~\ref{hyp:global}, Theorem~\ref{thm:main} included.

Finally, the exact constant $c^{**}=0.382885260\ldots$ is not computed as a fixed point; the
formal statement is \texttt{refined\_ceiling\_numeric}, $c\le3829/10000$. Its proof is the
paper's: \texttt{fixed\_point\_form} at $x=0.6909$ gives, after clearing the denominator,
$x\,\ent(x)\le(1-c)-(1-\ent(x^2))/2$, which is contradicted by $c>0.3829$ together with the
rigorous bounds
\[0.892131\le\ent(0.6909)\le0.8921319,\qquad 0.9985182\le\ent(0.6909^2)\le0.998519\]
(the lemmas \texttt{hx\_lower}, \texttt{hx\_upper}, \texttt{hx2\_lower},
\texttt{hx2\_upper} of \texttt{Refined.lean}, the last two stated at
$0.6909^2=0.47734281$; the true values are $0.89213183\ldots$ and $0.99851828\ldots$). Those
bounds come from Taylor bounds on $\log(1-t)$ with Mathlib's explicit remainder
\texttt{Real.abs\_log\_sub\_add\_sum\_range\_le} ($16$, $16$ and $6$ terms for
$\log0.6909$, $\log0.6182$ and $\log1.04531438$ respectively) together with Mathlib's
\texttt{Real.log\_two\_gt\_d9} and \texttt{Real.log\_two\_lt\_d9}. Since
$0.3829<c_{\mathrm{ceil}}=0.3830993\ldots$, the formal statement already shows that the
refined ceiling lies strictly below the product ceiling; the margin between $0.3829$ and
$c^{**}$ is $\approx1.5\cdot10^{-5}$.

\subsection{The certificate}\label{app:lean:check}
The development builds with Lean~4.23.0 (\texttt{leanprover/lean4:v4.23.0}) against Mathlib
pinned at tag \texttt{v4.23.0}, commit \texttt{37df177aaa770670452312393d4e84aaad56e7b6}. The
gate script \texttt{check.sh} (Appendix~\ref{app:src} lists the library sources; the script
itself is in the ancillary directory) runs \texttt{lake build}, scans the build output and
the sources for \texttt{sorry}, and runs \texttt{lake env lean AxiomCheck.lean}, requiring
that every checked declaration report no axioms beyond \texttt{propext},
\texttt{Classical.choice} and \texttt{Quot.sound}; it exits $0$ only if all three pass. It
does pass, with the \texttt{\#print axioms} output

\begin{Verbatim}[fontsize=\scriptsize,formatcom=\leanmono,samepage=false]
'UnionClosedCeiling.product_ceiling' depends on axioms: [propext, Classical.choice, Quot.sound]
'UnionClosedCeiling.hiding_bound' depends on axioms: [propext, Classical.choice, Quot.sound]
'UnionClosedCeiling.diagonal_bound' depends on axioms: [propext, Classical.choice, Quot.sound]
'UnionClosedCeiling.fixed_point_form' depends on axioms: [propext, Classical.choice, Quot.sound]
'UnionClosedCeiling.refined_ceiling_numeric' depends on axioms: [propext, Classical.choice, Quot.sound]
'UnionClosedCeiling.containsProduct_iidClass' depends on axioms: [propext, Classical.choice, Quot.sound]
'UnionClosedCeiling.containsProduct_allCouplings' depends on axioms: [propext, Classical.choice, Quot.sound]
'UnionClosedCeiling.containsProduct_mixtureOfProducts' depends on axioms: [propext, Classical.choice, Quot.sound]
'UnionClosedCeiling.admitsHiding_allCouplings' depends on axioms: [propext, Classical.choice, Quot.sound]
'UnionClosedCeiling.admitsHiding_mixtureOfProducts' depends on axioms: [propext, Classical.choice, Quot.sound]
\end{Verbatim}

The same gate runs in GitHub Actions on a fresh Ubuntu clone, where it takes 2\,min\,20\,s.
The code is the ancillary directory \texttt{anc/lean/} of the arXiv package and the directory
\texttt{lean/} of the repository \url{https://github.com/moffatstudio/union-closed-constant};
from a fresh clone of either,
\[\texttt{lake exe cache get \&\& lake build \&\& bash check.sh}\]
reproduces the check.


\section{Lean source}\label{app:src}
The complete source of the library described in Appendix~\ref{app:lean}, in build order. The
same files are the ancillary material of this submission; the listings below are copies of
them, character for character. (The text layer of the PDF is not a faithful copy: the listing font maps the ASCII hyphen and the middle dot to typographic variants, so the source should be taken from \texttt{anc/lean/}, not copied out of the PDF.)

\subsection{\texttt{UnionClosedCeiling/Entropy.lean}}
\begin{Verbatim}[fontsize=\scriptsize,formatcom=\leanmono,samepage=false]
import Mathlib.Analysis.SpecialFunctions.BinaryEntropy

/-!
# Binary entropy in bits

The paper writes `ent` for the binary entropy function measured in **bits**,
`ent(p) = -p log₂ p - (1-p) log₂ (1-p)`, normalised so that `ent(1/2) = 1`.
Mathlib's `Real.binEntropy` is the same function in *nats*, so we set
`h p = Real.binEntropy p / Real.log 2`.

Facts proved here (all of them are used later):
`h_zero`, `h_one`, `h_nonneg`, `h_pos`, `h_le_one`.
-/

namespace UnionClosedCeiling

open Real

/-- Binary entropy in bits: `h p = -p log₂ p - (1-p) log₂ (1-p)`.
This is the paper's `ent`. -/
noncomputable def h (p : ℝ) : ℝ := Real.binEntropy p / Real.log 2

lemma log_two_pos : (0:ℝ) < Real.log 2 := Real.log_pos (by norm_num)

@[simp] lemma h_zero : h 0 = 0 := by simp [h]

@[simp] lemma h_one : h 1 = 0 := by simp [h]

@[simp] lemma h_half : h (1/2) = 1 := by
  have : (1:ℝ)/2 = (2:ℝ)⁻¹ := by norm_num
  simp [h, this, Real.binEntropy_two_inv, div_self (ne_of_gt log_two_pos)]

/-- `h` is nonnegative on `[0,1]`. -/
lemma h_nonneg {p : ℝ} (hp₀ : 0 ≤ p) (hp₁ : p ≤ 1) : 0 ≤ h p :=
  div_nonneg (Real.binEntropy_nonneg hp₀ hp₁) (le_of_lt log_two_pos)

/-- `h` is strictly positive on `(0,1)`. -/
lemma h_pos {p : ℝ} (hp₀ : 0 < p) (hp₁ : p < 1) : 0 < h p :=
  div_pos (Real.binEntropy_pos hp₀ hp₁) log_two_pos

/-- `h p ≤ 1` for every real `p`: the paper's `ent(·) ≤ 1`. -/
lemma h_le_one (p : ℝ) : h p ≤ 1 := by
  unfold h
  rw [div_le_one log_two_pos]
  exact Real.binEntropy_le_log_two

end UnionClosedCeiling
\end{Verbatim}

\subsection{\texttt{UnionClosedCeiling/Framework.lean}}
\begin{Verbatim}[fontsize=\scriptsize,formatcom=\leanmono,samepage=false]
import UnionClosedCeiling.Entropy

/-!
# The single-letter framework (finite model)

This file sets up a faithful **finite** model of the framework of Section 2.2 of the paper:
laws on `[0,1]` with finite support, couplings of such a law with itself, protocols
(represented by their value `Prot_{x,y}(0,0)`, which is all Section 3 uses about them),
classes of couplings, and the certificate inequality (cert).

## Modelling choices

* A law `μ` is given by atoms `a : Fin n → ℝ` in `[0,1]` and weights `m : Fin n → ℝ`
  summing to `1`.  All the adversarial laws used in Section 3 (two-point laws and the
  hiding law of Definition 3.2) are of this form, so restricting to finitely supported
  laws only *weakens* the certificate hypothesis, i.e. strengthens the theorems.
* A protocol is modelled by the function `Prot : ℝ → ℝ → ℝ`, `Prot x y = Prot_{x,y}(0,0)`,
  subject to the Fréchet bounds `max 0 (x+y-1) ≤ Prot x y ≤ min x y`.  Section 3 uses
  nothing else about protocols.
* **`Certifies` is implied by the paper's inequality (cert).**  (cert) reads
  `∑ₖ wₖ inf_{P ∈ Cₖ(μ)} 𝔼_P[ent(Prot⁽ᵏ⁾(X,Y))] ≥ C 𝔼_μ[ent X]`; since the infimum is at
  most the value at any particular admissible family `(P k)`, (cert) implies the
  statement `Certifies` below for every choice of admissible couplings.  Hence every
  theorem proved from `Certifies` applies verbatim to the paper's certificate.
-/

namespace UnionClosedCeiling

open Finset

variable {n : ℕ}

/-- `a` are atoms in `[0,1]` and `m` are probability weights: a finitely supported law on `[0,1]`. -/
def IsLaw (a m : Fin n → ℝ) : Prop :=
  (∀ i, a i ∈ Set.Icc (0:ℝ) 1) ∧ (∀ i, 0 ≤ m i) ∧ ∑ i, m i = 1

/-- The mean `𝔼_μ[X]` of the law. -/
noncomputable def mean (a m : Fin n → ℝ) : ℝ := ∑ i, m i * a i

/-- The entropy `𝔼_μ[ent X]` of the law, in bits. -/
noncomputable def entropy (a m : Fin n → ℝ) : ℝ := ∑ i, m i * h (a i)

/-- `P` is a coupling of the law with weights `m` with itself. -/
def IsCoupling (m : Fin n → ℝ) (P : Fin n → Fin n → ℝ) : Prop :=
  (∀ i j, 0 ≤ P i j) ∧ (∀ i, ∑ j, P i j = m i) ∧ (∀ j, ∑ i, P i j = m j)

/-- `𝔼_{(X,Y) ∼ P}[f X Y]`. -/
noncomputable def expect (a : Fin n → ℝ) (P : Fin n → Fin n → ℝ) (f : ℝ → ℝ → ℝ) : ℝ :=
  ∑ i, ∑ j, P i j * f (a i) (a j)

/-- The product coupling `μ ⊗ μ`. -/
def productCoupling (m : Fin n → ℝ) : Fin n → Fin n → ℝ := fun i j => m i * m j

lemma isCoupling_productCoupling {m : Fin n → ℝ} (hm : ∀ i, 0 ≤ m i) (hs : ∑ i, m i = 1) :
    IsCoupling m (productCoupling m) := by
  refine ⟨fun i j => mul_nonneg (hm i) (hm j), fun i => ?_, fun j => ?_⟩
  · simp [productCoupling, ← Finset.mul_sum, hs]
  · simp [productCoupling, ← Finset.sum_mul, hs]

/-- A protocol, represented by `Prot x y = Prot_{x,y}(0,0)`, constrained by the Fréchet bounds. -/
def IsProtocol (Prot : ℝ → ℝ → ℝ) : Prop :=
  ∀ x y, x ∈ Set.Icc (0:ℝ) 1 → y ∈ Set.Icc (0:ℝ) 1 →
    max 0 (x + y - 1) ≤ Prot x y ∧ Prot x y ≤ min x y

section ProtocolFacts

variable {Prot : ℝ → ℝ → ℝ} (hProt : IsProtocol Prot)
include hProt

/-- "Pairs containing the atom `0` carry no entropy": `Prot x 0 = 0`. -/
lemma protocol_right_zero {x : ℝ} (hx : x ∈ Set.Icc (0:ℝ) 1) : Prot x 0 = 0 := by
  obtain ⟨hl, hu⟩ := hProt x 0 hx ⟨le_refl 0, by norm_num⟩
  have h1 : (0:ℝ) ≤ Prot x 0 := le_trans (le_max_left _ _) hl
  have h2 : Prot x 0 ≤ 0 := le_trans hu (min_le_right _ _)
  linarith

lemma protocol_left_zero {y : ℝ} (hy : y ∈ Set.Icc (0:ℝ) 1) : Prot 0 y = 0 := by
  obtain ⟨hl, hu⟩ := hProt 0 y ⟨le_refl 0, by norm_num⟩ hy
  have h1 : (0:ℝ) ≤ Prot 0 y := le_trans (le_max_left _ _) hl
  have h2 : Prot 0 y ≤ 0 := le_trans hu (min_le_left _ _)
  linarith

/-- On the pair `(1,1)` both bits are surely `0`: `Prot 1 1 = 1`. -/
lemma protocol_one_one : Prot 1 1 = 1 := by
  obtain ⟨hl, hu⟩ := hProt 1 1 ⟨by norm_num, le_refl 1⟩ ⟨by norm_num, le_refl 1⟩
  have h1 : (1:ℝ) ≤ Prot 1 1 := by
    have : (1:ℝ) ≤ max 0 (1 + 1 - 1) := by norm_num
    linarith [le_trans this hl]
  have h2 : Prot 1 1 ≤ 1 := le_trans hu (min_le_left _ _)
  linarith

/-- On the pair `(1,y)` one bit is surely `0`, so the union bit equals the other bit. -/
lemma protocol_one_left {y : ℝ} (hy : y ∈ Set.Icc (0:ℝ) 1) : Prot 1 y = y := by
  obtain ⟨hl, hu⟩ := hProt 1 y ⟨by norm_num, le_refl 1⟩ hy
  have h1 : y ≤ Prot 1 y := by
    have : y ≤ max 0 (1 + y - 1) := by
      have := le_max_right (0:ℝ) (1 + y - 1); linarith
    linarith [le_trans this hl]
  have h2 : Prot 1 y ≤ y := le_trans hu (min_le_right _ _)
  linarith

lemma protocol_one_right {x : ℝ} (hx : x ∈ Set.Icc (0:ℝ) 1) : Prot x 1 = x := by
  obtain ⟨hl, hu⟩ := hProt x 1 hx ⟨by norm_num, le_refl 1⟩
  have h1 : x ≤ Prot x 1 := by
    have : x ≤ max 0 (x + 1 - 1) := by
      have := le_max_right (0:ℝ) (x + 1 - 1); linarith
    linarith [le_trans this hl]
  have h2 : Prot x 1 ≤ x := le_trans hu (min_le_left _ _)
  linarith

end ProtocolFacts

variable {K : ℕ}

/-- The certificate inequality (cert) of Proposition 2.2, in the form we use:
for every finitely supported law of mean at least `1 - c` and every admissible family
of couplings `P k` (one in each class `C k`), the weighted average of the protocol
entropies dominates `Cc` times the entropy of the law.

This is **implied by** the paper's (cert), whose left-hand side is an infimum over the
class, hence at most the value at the chosen `P k`. -/
def Certifies (Prot : Fin K → ℝ → ℝ → ℝ) (w : Fin K → ℝ)
    (C : Fin K → ∀ n : ℕ, (Fin n → ℝ) → (Fin n → ℝ) → (Fin n → Fin n → ℝ) → Prop)
    (Cc c : ℝ) : Prop :=
  ∀ (n : ℕ) (a m : Fin n → ℝ), IsLaw a m → 1 - c ≤ mean a m →
    ∀ P : Fin K → (Fin n → Fin n → ℝ),
      (∀ k, IsCoupling m (P k) ∧ C k n a m (P k)) →
      Cc * entropy a m ≤ ∑ k, w k * expect a (P k) (fun x y => h (Prot k x y))

/-- The class `C` contains the product law `μ ⊗ μ` for every law `μ`. -/
def ContainsProduct (C : ∀ n : ℕ, (Fin n → ℝ) → (Fin n → ℝ) → (Fin n → Fin n → ℝ) → Prop) :
    Prop :=
  ∀ (n : ℕ) (a m : Fin n → ℝ), IsLaw a m → C n a m (productCoupling m)

/-! ### The two-point adversary `p δ_x + (1-p) δ_0` -/

/-- Atoms of the two-point law: `x` and `0`. -/
def twoAtoms (x : ℝ) : Fin 2 → ℝ := ![x, 0]

/-- Weights of the two-point law: `p` and `1-p`. -/
def twoWeights (p : ℝ) : Fin 2 → ℝ := ![p, 1 - p]

lemma isLaw_two {x p : ℝ} (hx : x ∈ Set.Icc (0:ℝ) 1) (hp₀ : 0 ≤ p) (hp₁ : p ≤ 1) :
    IsLaw (twoAtoms x) (twoWeights p) := by
  refine ⟨?_, ?_, ?_⟩
  · intro i; fin_cases i
    · exact hx
    · simp [twoAtoms, Set.mem_Icc]
  · intro i; fin_cases i
    · exact hp₀
    · show (0:ℝ) ≤ 1 - p
      linarith
  · simp [twoWeights, Fin.sum_univ_two]

lemma mean_two {x p : ℝ} : mean (twoAtoms x) (twoWeights p) = p * x := by
  simp [mean, twoAtoms, twoWeights, Fin.sum_univ_two]

lemma entropy_two {x p : ℝ} : entropy (twoAtoms x) (twoWeights p) = p * h x := by
  simp [entropy, twoAtoms, twoWeights, Fin.sum_univ_two]

/-- Under the product coupling, the two-point law puts all the entropy on the diagonal
cell `(x,x)`: every other cell contains the atom `0`. -/
lemma expect_two_product {Prot : ℝ → ℝ → ℝ} (hProt : IsProtocol Prot) {x p : ℝ}
    (hx : x ∈ Set.Icc (0:ℝ) 1) :
    expect (twoAtoms x) (productCoupling (twoWeights p)) (fun u v => h (Prot u v))
      = p * p * h (Prot x x) := by
  have h0 : (0:ℝ) ∈ Set.Icc (0:ℝ) 1 := ⟨le_refl 0, by norm_num⟩
  simp [expect, productCoupling, twoAtoms, twoWeights, Fin.sum_univ_two,
    protocol_right_zero hProt hx, protocol_left_zero hProt hx, protocol_right_zero hProt h0]

/-! ### The hiding adversary `μ_{q,y,δ}` (Definition 3.2)

`μ_{q,y,δ} = q δ₁ + (1-q)[(1-δ) δ₀ + δ δ_y]`, with index `0 ↦ 1`, `1 ↦ 0`, `2 ↦ y`. -/

/-- Atoms of the hiding law: `1`, `0`, `y`. -/
def hidingAtoms (y : ℝ) : Fin 3 → ℝ := ![1, 0, y]

/-- Weights of the hiding law: `q`, `(1-q)(1-δ)`, `(1-q)δ`. -/
def hidingWeights (q δ : ℝ) : Fin 3 → ℝ := ![q, (1-q)*(1-δ), (1-q)*δ]

lemma isLaw_hiding {q y δ : ℝ} (hq₀ : 0 ≤ q) (hq₁ : q ≤ 1) (hδ₀ : 0 ≤ δ) (hδ₁ : δ ≤ 1)
    (hy : y ∈ Set.Icc (0:ℝ) 1) : IsLaw (hidingAtoms y) (hidingWeights q δ) := by
  refine ⟨?_, ?_, ?_⟩
  · intro i; fin_cases i
    · simp [hidingAtoms, Set.mem_Icc]
    · simp [hidingAtoms, Set.mem_Icc]
    · exact hy
  · intro i; fin_cases i
    · exact hq₀
    · exact mul_nonneg (by linarith) (by linarith)
    · exact mul_nonneg (by linarith) hδ₀
  · simp [hidingWeights, Fin.sum_univ_three]; ring

lemma mean_hiding {q y δ : ℝ} :
    mean (hidingAtoms y) (hidingWeights q δ) = q + (1-q)*δ*y := by
  simp [mean, hidingAtoms, hidingWeights, Fin.sum_univ_three]

lemma entropy_hiding {q y δ : ℝ} :
    entropy (hidingAtoms y) (hidingWeights q δ) = (1-q)*δ * h y := by
  simp [entropy, hidingAtoms, hidingWeights, Fin.sum_univ_three]

/-- The class `C` admits hiding (Definition 3.2): for all `q, y ∈ (0,1)` and all
sufficiently small `δ > 0` it contains a coupling of `μ_{q,y,δ}` with itself that puts
no mass on `(y,1)` and `(1,y)` and mass at most `δ²` on `(y,y)`. -/
def AdmitsHiding (C : ∀ n : ℕ, (Fin n → ℝ) → (Fin n → ℝ) → (Fin n → Fin n → ℝ) → Prop) :
    Prop :=
  ∀ q y : ℝ, q ∈ Set.Ioo (0:ℝ) 1 → y ∈ Set.Ioo (0:ℝ) 1 →
    ∃ δ₀ > (0:ℝ), ∀ δ : ℝ, 0 < δ → δ < δ₀ →
      ∃ P : Fin 3 → Fin 3 → ℝ, IsCoupling (hidingWeights q δ) P ∧
        C 3 (hidingAtoms y) (hidingWeights q δ) P ∧
        P 2 0 = 0 ∧ P 0 2 = 0 ∧ P 2 2 ≤ δ^2

/-- Value of the hiding law under a hiding coupling: at most `δ²`. -/
lemma expect_hiding_le {Prot : ℝ → ℝ → ℝ} (hProt : IsProtocol Prot) {q y δ : ℝ}
    (hy : y ∈ Set.Icc (0:ℝ) 1) {P : Fin 3 → Fin 3 → ℝ}
    (hP : IsCoupling (hidingWeights q δ) P)
    (h20 : P 2 0 = 0) (h02 : P 0 2 = 0) (h22 : P 2 2 ≤ δ^2) :
    expect (hidingAtoms y) P (fun u v => h (Prot u v)) ≤ δ^2 := by
  have h0 : (0:ℝ) ∈ Set.Icc (0:ℝ) 1 := ⟨le_refl 0, by norm_num⟩
  have h1 : (1:ℝ) ∈ Set.Icc (0:ℝ) 1 := ⟨by norm_num, le_refl 1⟩
  have hnn := hP.1
  have hkey : expect (hidingAtoms y) P (fun u v => h (Prot u v)) = P 2 2 * h (Prot y y) := by
    simp [expect, hidingAtoms, Fin.sum_univ_three, h20, h02,
      protocol_right_zero hProt h1, protocol_right_zero hProt hy, protocol_right_zero hProt h0,
      protocol_left_zero hProt h1, protocol_left_zero hProt hy,
      protocol_one_one hProt]
  rw [hkey]
  calc P 2 2 * h (Prot y y) ≤ P 2 2 * 1 :=
        mul_le_mul_of_nonneg_left (h_le_one _) (hnn 2 2)
    _ = P 2 2 := by ring
    _ ≤ δ^2 := h22

/-- Value of the hiding law under the **product** coupling, for the i.i.d. protocol
`Prot x y = x y`: exactly `2q(1-q)δ ent(y) + (1-q)²δ² ent(y²)`. -/
lemma expect_hiding_product_iid {q y δ : ℝ} :
    expect (hidingAtoms y) (productCoupling (hidingWeights q δ)) (fun u v => h (u * v))
      = 2*q*((1-q)*δ) * h y + ((1-q)*δ)^2 * h (y^2) := by
  have h1y : (1:ℝ) * y = y := one_mul y
  have hy1 : y * (1:ℝ) = y := mul_one y
  have hyy : y * y = y^2 := (sq y).symm
  simp [expect, productCoupling, hidingAtoms, hidingWeights, Fin.sum_univ_three,
    h1y, hyy]
  ring

end UnionClosedCeiling
\end{Verbatim}

\subsection{\texttt{UnionClosedCeiling/Ceiling.lean}}
\begin{Verbatim}[fontsize=\scriptsize,formatcom=\leanmono,samepage=false]
import UnionClosedCeiling.Framework

/-!
# Theorem 3.1 — the product ceiling

> **Theorem (product ceiling).**  Suppose every class `C_k(μ)` in (cert) contains the
> product law `μ ⊗ μ` for every `μ`.  Then (cert) fails for every
> `c > c_ceil := 1 - ent(1/√2)/√2 = 0.3830993…`.

Formal contrapositive: if the certificate holds for `c` with `C > 1` (we only use
`C ≥ 1`), then `c ≤ 1 - ent(1/√2)/√2`.

The proof only uses `ent ≤ 1` on the diagonal cell — **not** `ent(x*²) = 1`.
-/

namespace UnionClosedCeiling

open Finset

/-- The ceiling constant `c_ceil = 1 - ent(1/√2)/√2 = 0.3830993…`. -/
noncomputable def cCeil : ℝ := 1 - h (1 / Real.sqrt 2) / Real.sqrt 2

lemma one_lt_sqrt_two : (1:ℝ) < Real.sqrt 2 := by
  nlinarith [Real.sq_sqrt (by norm_num : (0:ℝ) ≤ 2), Real.sqrt_nonneg 2]

lemma sqrt_two_pos : (0:ℝ) < Real.sqrt 2 := lt_trans one_pos one_lt_sqrt_two

/-- `x* = 1/√2` lies in `(0,1)`. -/
lemma xstar_mem : (1 / Real.sqrt 2) ∈ Set.Ioo (0:ℝ) 1 := by
  constructor
  · positivity
  · rw [div_lt_one sqrt_two_pos]; exact one_lt_sqrt_two

/-- **Theorem 3.1 (product ceiling).**  If every class contains the product law and the
certificate (cert) holds for `c` with constant `Cc ≥ 1`, then `c ≤ 1 - ent(1/√2)/√2`. -/
theorem product_ceiling {K : ℕ} {Prot : Fin K → ℝ → ℝ → ℝ} {w : Fin K → ℝ}
    {C : Fin K → ∀ n : ℕ, (Fin n → ℝ) → (Fin n → ℝ) → (Fin n → Fin n → ℝ) → Prop}
    {Cc c : ℝ}
    (hProt : ∀ k, IsProtocol (Prot k)) (hw : ∀ k, 0 ≤ w k) (hw1 : ∑ k, w k = 1)
    (hC : ∀ k, ContainsProduct (C k)) (hCc : 1 ≤ Cc)
    (hcert : Certifies Prot w C Cc c) :
    c ≤ 1 - h (1 / Real.sqrt 2) / Real.sqrt 2 := by
  by_contra hcon
  push_neg at hcon
  set x := 1 / Real.sqrt 2 with hxdef
  obtain ⟨hx0, hx1⟩ := xstar_mem
  have hxIcc : x ∈ Set.Icc (0:ℝ) 1 := ⟨le_of_lt hx0, le_of_lt hx1⟩
  have hhx : 0 < h x := h_pos hx0 hx1
  -- `(1-c)·√2 < ent(x*)`
  have hkey : (1 - c) * Real.sqrt 2 < h x := by
    have : 1 - h x / Real.sqrt 2 < c := hcon
    have h2 : 1 - c < h x / Real.sqrt 2 := by linarith
    calc (1 - c) * Real.sqrt 2 < (h x / Real.sqrt 2) * Real.sqrt 2 := by
          exact mul_lt_mul_of_pos_right h2 sqrt_two_pos
      _ = h x := by field_simp
  -- choose `p` strictly between `max 0 ((1-c)√2)` and `ent(x*)`
  set L : ℝ := max 0 ((1 - c) * Real.sqrt 2) with hLdef
  have hL : L < h x := max_lt hhx hkey
  have hL0 : 0 ≤ L := le_max_left _ _
  set p : ℝ := (L + h x) / 2 with hpdef
  have hpL : L < p := by simp only [hpdef]; linarith
  have hph : p < h x := by simp only [hpdef]; linarith
  have hp0 : 0 < p := lt_of_le_of_lt hL0 hpL
  have hp1 : p ≤ 1 := le_of_lt (lt_of_lt_of_le hph (h_le_one x))
  -- the two-point law `p δ_{x*} + (1-p) δ_0` has mean `p x* ≥ 1 - c`
  have hlaw : IsLaw (twoAtoms x) (twoWeights p) := isLaw_two hxIcc (le_of_lt hp0) hp1
  have hmean : 1 - c ≤ mean (twoAtoms x) (twoWeights p) := by
    rw [mean_two]
    have h1 : (1 - c) * Real.sqrt 2 < p := lt_of_le_of_lt (le_max_right _ _) hpL
    have h2 : 1 - c ≤ p / Real.sqrt 2 := by
      rw [le_div_iff₀ sqrt_two_pos]; linarith
    calc 1 - c ≤ p / Real.sqrt 2 := h2
      _ = p * x := by rw [hxdef]; ring
  -- present the product coupling to every class
  have hcoup : ∀ k, IsCoupling (twoWeights p) (productCoupling (twoWeights p)) ∧
      C k 2 (twoAtoms x) (twoWeights p) (productCoupling (twoWeights p)) := by
    intro k
    exact ⟨isCoupling_productCoupling hlaw.2.1 hlaw.2.2, hC k 2 _ _ hlaw⟩
  have hmain := hcert 2 (twoAtoms x) (twoWeights p) hlaw hmean
    (fun _ => productCoupling (twoWeights p)) hcoup
  rw [entropy_two] at hmain
  -- every term is at most `w k * p²`
  have hsum : ∑ k, w k * expect (twoAtoms x) (productCoupling (twoWeights p))
      (fun u v => h (Prot k u v)) ≤ p * p := by
    have hterm : ∀ k ∈ (univ : Finset (Fin K)),
        w k * expect (twoAtoms x) (productCoupling (twoWeights p))
          (fun u v => h (Prot k u v)) ≤ w k * (p * p) := by
      intro k _
      rw [expect_two_product (hProt k) hxIcc]
      refine mul_le_mul_of_nonneg_left ?_ (hw k)
      have := h_le_one (Prot k x x)
      nlinarith [mul_pos hp0 hp0]
    calc ∑ k, w k * expect (twoAtoms x) (productCoupling (twoWeights p))
            (fun u v => h (Prot k u v))
        ≤ ∑ k, w k * (p * p) := Finset.sum_le_sum hterm
      _ = p * p := by rw [← Finset.sum_mul, hw1, one_mul]
  -- hence `p ent(x*) ≤ Cc p ent(x*) ≤ p²`, so `ent(x*) ≤ p`, contradiction
  have h1 : p * h x ≤ Cc * (p * h x) := by nlinarith [mul_nonneg (le_of_lt hp0) (le_of_lt hhx)]
  have h2 : p * h x ≤ p * p := le_trans h1 (le_trans hmain hsum)
  have h3 : h x ≤ p := le_of_mul_le_mul_left (by linarith) hp0
  linarith
\end{Verbatim}

\subsection{\texttt{UnionClosedCeiling/Refined.lean}}
\begin{Verbatim}[fontsize=\scriptsize,formatcom=\leanmono,samepage=false]
import UnionClosedCeiling.Framework
import Mathlib.Analysis.SpecialFunctions.Log.Deriv
import Mathlib.Analysis.Complex.ExponentialBounds

/-!
# Theorem 3.4 — the refined ceiling

> **Theorem (refined ceiling).**  Suppose (cert) uses the i.i.d. protocol with its
> singleton class `{μ ⊗ μ}` and weight `w ∈ [0,1]`, and that every other class contains
> the product law and admits hiding.  Then
> `2w(1-c) ≥ 1` and `c ≤ 1 - x ent(x)/(w ent(x²) + 1 - w)` for every `x` with
> `ent(x) ≤ w ent(x²) + 1 - w`.  Consequently `c ≤ c** = 0.382885260…`.

Formalised here: the two inequalities (H) = `hiding_bound` and (D) = `diagonal_bound`,
the exact fixed-point form `fixed_point_form` they combine into, and the numerical
corollary `refined_ceiling_numeric` (`c ≤ 0.3829`, which already lies below the product
ceiling `c_ceil = 0.3830993…`).

Protocol `0` is the i.i.d. one.  The hypothesis on its class is only
`ContainsProduct (C 0)`, which is weaker than the paper's singleton class, so the
theorems are correspondingly stronger.
-/

namespace UnionClosedCeiling

open Finset

/-! ### Rigorous numerical bounds on the binary entropy at `x = 0.6909`

All bounds come from the Taylor series of `log (1-t)` with the explicit remainder
`Real.abs_log_sub_add_sum_range_le`, together with Mathlib's `log 2` bounds. -/

lemma binEntropy_eq' (p : ℝ) :
    Real.binEntropy p = -(p * Real.log p) - (1 - p) * Real.log (1 - p) := by
  simp [Real.binEntropy, Real.log_inv]; ring

lemma h_eq (p : ℝ) : h p = (-(p * Real.log p) - (1 - p) * Real.log (1 - p)) / Real.log 2 := by
  rw [h, binEntropy_eq']

/-- `log 0.6909 ∈ [-0.36976019, -0.36976018]`. -/
lemma log_6909 : (-0.36976019 : ℝ) ≤ Real.log (6909/10000) ∧
    Real.log (6909/10000) ≤ -0.36976018 := by
  have ht : |(3091/10000 : ℝ)| < 1 := by rw [abs_lt]; norm_num
  have key := Real.abs_log_sub_add_sum_range_le ht 16
  rw [abs_le] at key
  obtain ⟨k1, k2⟩ := key
  have h1 : (1:ℝ) - 3091/10000 = 6909/10000 := by norm_num
  rw [h1] at k1 k2
  rw [abs_of_pos (by norm_num : (0:ℝ) < 3091/10000)] at k1 k2
  have hA : (∑ i ∈ Finset.range 16, ((3091:ℝ)/10000) ^ (i + 1) / (i + 1))
      + ((3091:ℝ)/10000) ^ (16 + 1) / (1 - 3091/10000) ≤ 0.36976019 := by
    norm_num [Finset.sum_range_succ]
  have hB : (0.36976018 : ℝ) ≤ (∑ i ∈ Finset.range 16, ((3091:ℝ)/10000) ^ (i + 1) / (i + 1))
      - ((3091:ℝ)/10000) ^ (16 + 1) / (1 - 3091/10000) := by
    norm_num [Finset.sum_range_succ]
  exact ⟨by linarith, by linarith⟩

/-- `log 0.6182 ∈ [-0.48094337, -0.48094311]`. -/
lemma log_6182 : (-0.48094337 : ℝ) ≤ Real.log (6182/10000) ∧
    Real.log (6182/10000) ≤ -0.48094311 := by
  have ht : |(3818/10000 : ℝ)| < 1 := by rw [abs_lt]; norm_num
  have key := Real.abs_log_sub_add_sum_range_le ht 16
  rw [abs_le] at key
  obtain ⟨k1, k2⟩ := key
  have h1 : (1:ℝ) - 3818/10000 = 6182/10000 := by norm_num
  rw [h1] at k1 k2
  rw [abs_of_pos (by norm_num : (0:ℝ) < 3818/10000)] at k1 k2
  have hA : (∑ i ∈ Finset.range 16, ((3818:ℝ)/10000) ^ (i + 1) / (i + 1))
      + ((3818:ℝ)/10000) ^ (16 + 1) / (1 - 3818/10000) ≤ 0.48094337 := by
    norm_num [Finset.sum_range_succ]
  have hB : (0.48094311 : ℝ) ≤ (∑ i ∈ Finset.range 16, ((3818:ℝ)/10000) ^ (i + 1) / (i + 1))
      - ((3818:ℝ)/10000) ^ (16 + 1) / (1 - 3818/10000) := by
    norm_num [Finset.sum_range_succ]
  exact ⟨by linarith, by linarith⟩

/-- `log 1.04531438 ∈ [0.04431768, 0.04431769]`. -/
lemma log_104531438 : (0.04431768 : ℝ) ≤ Real.log (104531438/100000000) ∧
    Real.log (104531438/100000000) ≤ 0.04431769 := by
  have ht : |(-4531438/100000000 : ℝ)| < 1 := by rw [abs_lt]; norm_num
  have key := Real.abs_log_sub_add_sum_range_le ht 6
  rw [abs_le] at key
  obtain ⟨k1, k2⟩ := key
  have h1 : (1:ℝ) - (-4531438/100000000) = 104531438/100000000 := by norm_num
  rw [h1] at k1 k2
  rw [abs_of_neg (by norm_num : (-4531438/100000000 : ℝ) < 0)] at k1 k2
  have hA : (∑ i ∈ Finset.range 6, ((-4531438:ℝ)/100000000) ^ (i + 1) / (i + 1))
      + (-((-4531438:ℝ)/100000000)) ^ (6 + 1) / (1 - -((-4531438:ℝ)/100000000))
      ≤ -0.04431768 := by
    norm_num [Finset.sum_range_succ]
  have hB : (-0.04431769 : ℝ) ≤
      (∑ i ∈ Finset.range 6, ((-4531438:ℝ)/100000000) ^ (i + 1) / (i + 1))
      - (-((-4531438:ℝ)/100000000)) ^ (6 + 1) / (1 - -((-4531438:ℝ)/100000000)) := by
    norm_num [Finset.sum_range_succ]
  exact ⟨by linarith, by linarith⟩

lemma log_3091 : Real.log (3091/10000) = Real.log (6182/10000) - Real.log 2 := by
  rw [show (3091:ℝ)/10000 = (6182/10000)/2 by norm_num,
    Real.log_div (by norm_num) (by norm_num)]

lemma log_52265719 :
    Real.log (52265719/100000000) = Real.log (104531438/100000000) - Real.log 2 := by
  rw [show (52265719:ℝ)/100000000 = (104531438/100000000)/2 by norm_num,
    Real.log_div (by norm_num) (by norm_num)]

lemma log_47734281 : Real.log (47734281/100000000) = 2 * Real.log (6909/10000) := by
  rw [show (47734281:ℝ)/100000000 = (6909/10000)^2 by norm_num, Real.log_pow]
  norm_num

lemma xsq_eq : ((6909:ℝ)/10000)^2 = 47734281/100000000 := by norm_num

/-- `ent(0.6909) ≥ 0.892131`. -/
lemma hx_lower : (892131/1000000 : ℝ) ≤ h (6909/10000) := by
  rw [h_eq, le_div_iff₀ log_two_pos]
  norm_num [log_3091]
  linarith [log_6909.1, log_6909.2, log_6182.1, log_6182.2,
    Real.log_two_gt_d9, Real.log_two_lt_d9]

/-- `ent(0.6909) ≤ 0.8921319`. -/
lemma hx_upper : h (6909/10000) ≤ (8921319/10000000 : ℝ) := by
  rw [h_eq, div_le_iff₀ log_two_pos]
  norm_num [log_3091]
  linarith [log_6909.1, log_6909.2, log_6182.1, log_6182.2,
    Real.log_two_gt_d9, Real.log_two_lt_d9]

/-- `ent(0.6909²) ≥ 0.9985182`. -/
lemma hx2_lower : (9985182/10000000 : ℝ) ≤ h (47734281/100000000) := by
  rw [h_eq, le_div_iff₀ log_two_pos]
  norm_num [log_52265719, log_47734281]
  linarith [log_6909.1, log_6909.2, log_104531438.1, log_104531438.2,
    Real.log_two_gt_d9, Real.log_two_lt_d9]

/-- `ent(0.6909²) ≤ 0.998519`. -/
lemma hx2_upper : h (47734281/100000000) ≤ (998519/1000000 : ℝ) := by
  rw [h_eq, div_le_iff₀ log_two_pos]
  norm_num [log_52265719, log_47734281]
  linarith [log_6909.1, log_6909.2, log_104531438.1, log_104531438.2,
    Real.log_two_gt_d9, Real.log_two_lt_d9]

/-! ### The two inequalities of Theorem 3.4

Throughout, protocols are indexed by `Fin (K+1)` (so there is at least one), index `0`
is the i.i.d. protocol `Π⁰_{x,y}(0,0) = xy`. -/

variable {K : ℕ} {Prot : Fin (K+1) → ℝ → ℝ → ℝ} {w : Fin (K+1) → ℝ}
  {C : Fin (K+1) → ∀ n : ℕ, (Fin n → ℝ) → (Fin n → ℝ) → (Fin n → Fin n → ℝ) → Prop}
  {Cc c : ℝ}

lemma weight_zero_le_one (hw : ∀ k, 0 ≤ w k) (hw1 : ∑ k, w k = 1) : w 0 ≤ 1 := by
  calc w 0 ≤ ∑ k, w k := Finset.single_le_sum (fun i _ => hw i) (Finset.mem_univ 0)
    _ = 1 := hw1
/-- **Inequality (H) of Theorem 3.4 (component hiding).**

> Dividing by `𝔼[ent X] = (1-q) δ ent(y)` and letting `δ → 0`, (cert) forces `2 w q ≥ C > 1`.
> Hence `2w(1-c) ≥ 1` is necessary.

The limit `δ → 0` is done by hand: `δ` is chosen below the hiding threshold of every class
and below `c (1 - 2 q w₀)`. -/
theorem hiding_bound
    (hProt : ∀ k, IsProtocol (Prot k)) (hiid : ∀ x y, Prot 0 x y = x * y)
    (hw : ∀ k, 0 ≤ w k) (hw1 : ∑ k, w k = 1)
    (hC0 : ContainsProduct (C 0)) (hChide : ∀ k, k ≠ 0 → AdmitsHiding (C k))
    (hCc : 1 ≤ Cc) (hc0 : 0 < c) (hc1 : c < 1)
    (hcert : Certifies Prot w C Cc c) :
    1 ≤ 2 * w 0 * (1 - c) := by
  by_contra hcon
  push_neg at hcon
  obtain ⟨q, hqdef⟩ : ∃ q : ℝ, q = 1 - c := ⟨1 - c, rfl⟩
  have hq0 : 0 < q := by rw [hqdef]; linarith
  have hq1 : q < 1 := by rw [hqdef]; linarith
  have hw0 : 0 ≤ w 0 := hw 0
  have hcon' : 2 * w 0 * q < 1 := by rw [hqdef]; linarith
  have hgap : 0 < 1 - 2 * q * w 0 := by nlinarith
  -- every non-i.i.d. class supplies a hiding coupling below its own threshold `d k`
  have hex : ∀ k : Fin (K+1), ∃ e : ℝ, 0 < e ∧ ∀ δ : ℝ, 0 < δ → δ < e → k ≠ 0 →
      ∃ P : Fin 3 → Fin 3 → ℝ, IsCoupling (hidingWeights q δ) P ∧
        C k 3 (hidingAtoms (1/2)) (hidingWeights q δ) P ∧
        P 2 0 = 0 ∧ P 0 2 = 0 ∧ P 2 2 ≤ δ^2 := by
    intro k
    by_cases hk : k = 0
    · exact ⟨1, one_pos, fun δ _ _ hk0 => absurd hk hk0⟩
    · obtain ⟨e, he, hprop⟩ := hChide k hk q (1/2) ⟨hq0, hq1⟩ ⟨by norm_num, by norm_num⟩
      exact ⟨e, he, fun δ hδ hδd _ => hprop δ hδ hδd⟩
  choose d hd hdprop using hex
  -- one `δ` small enough for all of them, and small enough to force the contradiction
  obtain ⟨δ, hδpos, hδd, hδsmall, hδ1⟩ :
      ∃ δ : ℝ, 0 < δ ∧ (∀ k, δ < d k) ∧ δ < c * (1 - 2*q*w 0) ∧ δ ≤ 1 := by
    obtain ⟨δ₀, hδ₀, hδ₀le⟩ : ∃ e : ℝ, 0 < e ∧ ∀ k, e ≤ d k :=
      ⟨Finset.univ.inf' Finset.univ_nonempty d,
        (Finset.lt_inf'_iff Finset.univ_nonempty).2 (fun i _ => hd i),
        fun k => Finset.inf'_le d (Finset.mem_univ k)⟩
    refine ⟨min (δ₀/2) (min (c * (1 - 2*q*w 0) / 2) (1/4)), ?_, ?_, ?_, ?_⟩
    · exact lt_min (by linarith) (lt_min (div_pos (mul_pos hc0 hgap) two_pos) (by norm_num))
    · intro k
      calc min (δ₀/2) (min (c * (1 - 2*q*w 0) / 2) (1/4)) ≤ δ₀/2 := min_le_left _ _
        _ < δ₀ := by linarith
        _ ≤ d k := hδ₀le k
    · calc min (δ₀/2) (min (c * (1 - 2*q*w 0) / 2) (1/4)) ≤ c * (1 - 2*q*w 0) / 2 :=
          le_trans (min_le_right _ _) (min_le_left _ _)
        _ < c * (1 - 2*q*w 0) := by nlinarith
    · calc min (δ₀/2) (min (c * (1 - 2*q*w 0) / 2) (1/4)) ≤ 1/4 :=
          le_trans (min_le_right _ _) (min_le_right _ _)
        _ ≤ 1 := by norm_num
  -- the hiding law `μ_{q,1/2,δ}`
  have hyIcc : (1/2 : ℝ) ∈ Set.Icc (0:ℝ) 1 := ⟨by norm_num, by norm_num⟩
  have hlaw : IsLaw (hidingAtoms (1/2)) (hidingWeights q δ) :=
    isLaw_hiding (le_of_lt hq0) (le_of_lt hq1) (le_of_lt hδpos) hδ1 hyIcc
  have hmean : 1 - c ≤ mean (hidingAtoms (1/2)) (hidingWeights q δ) := by
    rw [mean_hiding]
    have hpos : 0 < (1-q)*δ := mul_pos (by linarith) hδpos
    rw [hqdef] at hpos ⊢
    nlinarith
  -- product coupling for the i.i.d. protocol, a hiding coupling for the others
  have hPex : ∀ k : Fin (K+1), ∃ Pk : Fin 3 → Fin 3 → ℝ,
      IsCoupling (hidingWeights q δ) Pk ∧ C k 3 (hidingAtoms (1/2)) (hidingWeights q δ) Pk ∧
      (k ≠ 0 → expect (hidingAtoms (1/2)) Pk (fun u v => h (Prot k u v)) ≤ δ^2) ∧
      (k = 0 → Pk = productCoupling (hidingWeights q δ)) := by
    intro k
    by_cases hk : k = 0
    · subst hk
      exact ⟨productCoupling (hidingWeights q δ),
        isCoupling_productCoupling hlaw.2.1 hlaw.2.2, hC0 3 _ _ hlaw,
        fun hne => absurd rfl hne, fun _ => rfl⟩
    · obtain ⟨P, hP1, hP2, h20, h02, h22⟩ := hdprop k δ hδpos (hδd k) hk
      exact ⟨P, hP1, hP2,
        fun _ => expect_hiding_le (hProt k) hyIcc hP1 h20 h02 h22, fun he => absurd he hk⟩
  choose P hP1 hP2 hP3 hP4 using hPex
  have hmain := hcert 3 (hidingAtoms (1/2)) (hidingWeights q δ) hlaw hmean P
    (fun k => ⟨hP1 k, hP2 k⟩)
  rw [entropy_hiding] at hmain
  -- the i.i.d. term is exactly `2q(1-q)δ ent(y) + (1-q)²δ² ent(y²)`
  have hzero : expect (hidingAtoms (1/2)) (P 0) (fun u v => h (Prot 0 u v))
      = 2*q*((1-q)*δ) * h (1/2) + ((1-q)*δ)^2 * h ((1/2:ℝ)^2) := by
    rw [hP4 0 rfl]
    simp only [hiid]
    exact expect_hiding_product_iid
  -- every other term is at most `δ²`
  have hrest : ∑ k ∈ Finset.univ.erase (0 : Fin (K+1)),
      w k * expect (hidingAtoms (1/2)) (P k) (fun u v => h (Prot k u v))
      ≤ (1 - w 0) * δ^2 := by
    have hb : ∀ k ∈ Finset.univ.erase (0 : Fin (K+1)),
        w k * expect (hidingAtoms (1/2)) (P k) (fun u v => h (Prot k u v)) ≤ w k * δ^2 :=
      fun k hk => mul_le_mul_of_nonneg_left (hP3 k (Finset.ne_of_mem_erase hk)) (hw k)
    have hs : ∑ k ∈ Finset.univ.erase (0 : Fin (K+1)), w k = 1 - w 0 := by
      have := Finset.add_sum_erase Finset.univ w (Finset.mem_univ (0 : Fin (K+1)))
      rw [hw1] at this; linarith
    calc ∑ k ∈ Finset.univ.erase (0 : Fin (K+1)),
          w k * expect (hidingAtoms (1/2)) (P k) (fun u v => h (Prot k u v))
        ≤ ∑ k ∈ Finset.univ.erase (0 : Fin (K+1)), w k * δ^2 := Finset.sum_le_sum hb
      _ = (1 - w 0) * δ^2 := by rw [← Finset.sum_mul, hs]
  have hsplit : ∑ k, w k * expect (hidingAtoms (1/2)) (P k) (fun u v => h (Prot k u v))
      = w 0 * expect (hidingAtoms (1/2)) (P 0) (fun u v => h (Prot 0 u v))
        + ∑ k ∈ Finset.univ.erase (0 : Fin (K+1)),
            w k * expect (hidingAtoms (1/2)) (P k) (fun u v => h (Prot k u v)) :=
    (Finset.add_sum_erase Finset.univ _ (Finset.mem_univ (0 : Fin (K+1)))).symm
  rw [hsplit, hzero] at hmain
  have hhalf : h (1/2 : ℝ) = 1 := h_half
  have hquarter : h ((1/2:ℝ)^2) ≤ 1 := h_le_one _
  have hcq : (1 : ℝ) - q = c := by rw [hqdef]; ring
  rw [hhalf, hcq] at hmain
  -- arithmetic: `Cc c δ ≤ 2 q w₀ c δ + δ²`, contradicting `δ < c (1 - 2 q w₀)`
  have hsum2 : Cc * (c * δ * 1) ≤
      w 0 * (2*q*(c*δ) * 1 + (c*δ)^2 * h ((1/2:ℝ)^2)) + (1 - w 0) * δ^2 := by
    linarith [hmain, hrest]
  have hA : (c*δ)^2 * h ((1/2:ℝ)^2) ≤ δ^2 := by
    have e1 : (c*δ)^2 * h ((1/2:ℝ)^2) ≤ (c*δ)^2 := by
      nlinarith [mul_nonneg (sq_nonneg (c*δ)) (sub_nonneg.2 hquarter)]
    have e2 : (c*δ)^2 ≤ δ^2 := by
      nlinarith [mul_nonneg (sq_nonneg δ) (by nlinarith : (0:ℝ) ≤ 1 - c^2)]
    linarith
  have hstep : Cc * (c * δ) ≤ 2*q*w 0*(c*δ) + δ^2 := by
    linarith [hsum2, mul_nonneg hw0 (sub_nonneg.2 hA)]
  have h1 : c * δ ≤ 2*q*w 0*(c*δ) + δ^2 := by
    linarith [hstep, mul_nonneg (sub_nonneg.2 hCc) (le_of_lt (mul_pos hc0 hδpos))]
  have hfin : c ≤ 2*q*w 0*c + δ := by
    by_contra hno
    push_neg at hno
    have hprod : 0 < (c - 2*q*w 0*c - δ) * δ :=
      mul_pos (by linarith) hδpos
    linarith [h1, hprod]
  linarith [hfin, hδsmall]

/-- **Inequality (D) of Theorem 3.4 (the diagonal adversary).**

> Present the product law to every class at `μ = p δ_x + (1-p) δ_0`.  The i.i.d. term
> equals `p² ent(x²)` and every other term is at most `p²`, so (cert) forces
> `c ≤ 1 - x ent(x)/(w ent(x²) + 1 - w)`. -/
theorem diagonal_bound
    (hProt : ∀ k, IsProtocol (Prot k)) (hiid : ∀ x y, Prot 0 x y = x * y)
    (hw : ∀ k, 0 ≤ w k) (hw1 : ∑ k, w k = 1) (hC : ∀ k, ContainsProduct (C k))
    (hCc : 1 ≤ Cc) (hc1 : c < 1) (hcert : Certifies Prot w C Cc c)
    {x : ℝ} (hx : x ∈ Set.Ioc (0:ℝ) 1) (hcond : h x ≤ w 0 * h (x^2) + 1 - w 0) :
    c ≤ 1 - x * h x / (w 0 * h (x^2) + 1 - w 0) := by
  obtain ⟨hx0, hx1⟩ := hx
  rcases eq_or_lt_of_le hx1 with hxe | hxlt
  · subst hxe
    simp only [h_one, mul_zero, zero_div, sub_zero]
    linarith
  · have hxIcc : x ∈ Set.Icc (0:ℝ) 1 := ⟨le_of_lt hx0, hx1⟩
    obtain ⟨D, hDdef⟩ : ∃ D : ℝ, D = w 0 * h (x^2) + 1 - w 0 := ⟨_, rfl⟩
    have hhx : 0 < h x := h_pos hx0 hxlt
    have hcond' : h x ≤ D := by rw [hDdef]; exact hcond
    have hDpos : 0 < D := lt_of_lt_of_le hhx hcond'
    rw [← hDdef]
    by_contra hcon
    push_neg at hcon
    -- `p*` and a slightly smaller `p`
    obtain ⟨ps, hpsdef⟩ : ∃ ps : ℝ, ps = h x / D := ⟨_, rfl⟩
    have hps0 : 0 < ps := by rw [hpsdef]; exact div_pos hhx hDpos
    have hps1 : ps ≤ 1 := by rw [hpsdef, div_le_one hDpos]; exact hcond'
    have hxps : 1 - c < x * ps := by
      have hxx : x * h x / D = x * ps := by rw [hpsdef]; ring
      rw [hxx] at hcon; linarith
    obtain ⟨p, hplt, hp0, hpx⟩ : ∃ p : ℝ, p < ps ∧ 0 < p ∧ 1 - c ≤ p * x := by
      refine ⟨max ((1-c)/x) (ps/2), max_lt ?_ (by linarith), ?_, ?_⟩
      · rw [div_lt_iff₀ hx0]; nlinarith
      · exact lt_of_lt_of_le (by linarith : (0:ℝ) < ps/2) (le_max_right _ _)
      · have hle : (1-c)/x ≤ max ((1-c)/x) (ps/2) := le_max_left _ _
        rw [div_le_iff₀ hx0] at hle; linarith
    have hp1 : p ≤ 1 := le_of_lt (lt_of_lt_of_le hplt hps1)
    have hlaw : IsLaw (twoAtoms x) (twoWeights p) := isLaw_two hxIcc (le_of_lt hp0) hp1
    have hmean : 1 - c ≤ mean (twoAtoms x) (twoWeights p) := by rw [mean_two]; exact hpx
    have hmain := hcert 2 (twoAtoms x) (twoWeights p) hlaw hmean
      (fun _ => productCoupling (twoWeights p))
      (fun k => ⟨isCoupling_productCoupling hlaw.2.1 hlaw.2.2, hC k 2 _ _ hlaw⟩)
    rw [entropy_two] at hmain
    have hzero : expect (twoAtoms x) (productCoupling (twoWeights p))
        (fun u v => h (Prot 0 u v)) = p * p * h (x^2) := by
      have hxx : x * x = x^2 := by ring
      rw [expect_two_product (hProt 0) hxIcc, hiid, hxx]
    have hrest : ∑ k ∈ Finset.univ.erase (0 : Fin (K+1)), w k *
        expect (twoAtoms x) (productCoupling (twoWeights p)) (fun u v => h (Prot k u v))
        ≤ (1 - w 0) * (p * p) := by
      have hb : ∀ k ∈ Finset.univ.erase (0 : Fin (K+1)), w k *
          expect (twoAtoms x) (productCoupling (twoWeights p)) (fun u v => h (Prot k u v))
          ≤ w k * (p * p) := by
        intro k _
        rw [expect_two_product (hProt k) hxIcc]
        refine mul_le_mul_of_nonneg_left ?_ (hw k)
        nlinarith [h_le_one (Prot k x x), mul_pos hp0 hp0]
      have hs : ∑ k ∈ Finset.univ.erase (0 : Fin (K+1)), w k = 1 - w 0 := by
        have := Finset.add_sum_erase Finset.univ w (Finset.mem_univ (0 : Fin (K+1)))
        rw [hw1] at this; linarith
      calc ∑ k ∈ Finset.univ.erase (0 : Fin (K+1)), w k *
            expect (twoAtoms x) (productCoupling (twoWeights p)) (fun u v => h (Prot k u v))
          ≤ ∑ k ∈ Finset.univ.erase (0 : Fin (K+1)), w k * (p * p) := Finset.sum_le_sum hb
        _ = (1 - w 0) * (p * p) := by rw [← Finset.sum_mul, hs]
    have hsplit : ∑ k, w k *
        expect (twoAtoms x) (productCoupling (twoWeights p)) (fun u v => h (Prot k u v))
        = w 0 * expect (twoAtoms x) (productCoupling (twoWeights p))
            (fun u v => h (Prot 0 u v))
          + ∑ k ∈ Finset.univ.erase (0 : Fin (K+1)), w k *
              expect (twoAtoms x) (productCoupling (twoWeights p))
                (fun u v => h (Prot k u v)) :=
      (Finset.add_sum_erase Finset.univ _ (Finset.mem_univ (0 : Fin (K+1)))).symm
    rw [hsplit, hzero] at hmain
    have hfin : p * h x ≤ p * p * D := by
      have e1 : Cc * (p * h x) ≤ w 0 * (p * p * h (x^2)) + (1 - w 0) * (p * p) := by
        linarith [hmain, hrest]
      have e2 : w 0 * (p * p * h (x^2)) + (1 - w 0) * (p * p) = p * p * D := by
        rw [hDdef]; ring
      linarith [e1, e2, mul_nonneg (sub_nonneg.2 hCc)
        (mul_nonneg (le_of_lt hp0) (le_of_lt hhx))]
    have h3 : h x ≤ p * D := by
      by_contra hno
      push_neg at hno
      have hlt : p * (p * D) < p * h x := mul_lt_mul_of_pos_left hno hp0
      linarith [hfin, hlt]
    have h4 : ps ≤ p := by rw [hpsdef, div_le_iff₀ hDpos]; exact h3
    linarith

/-- **The fixed-point form of Theorem 3.4.**  Substituting the bound (H) `w₀ ≥ 1/(2(1-c))`
into (D), whose right-hand side is decreasing in `w₀`. -/
theorem fixed_point_form
    (hProt : ∀ k, IsProtocol (Prot k)) (hiid : ∀ x y, Prot 0 x y = x * y)
    (hw : ∀ k, 0 ≤ w k) (hw1 : ∑ k, w k = 1) (hC : ∀ k, ContainsProduct (C k))
    (hChide : ∀ k, k ≠ 0 → AdmitsHiding (C k))
    (hCc : 1 ≤ Cc) (hc0 : 0 < c) (hchalf : c ≤ 1/2) (hcert : Certifies Prot w C Cc c)
    {x : ℝ} (hx : x ∈ Set.Ioo (0:ℝ) 1) (hmono : h x ≤ h (x^2)) :
    c ≤ 1 - x * h x / (1 - (1 - h (x^2)) / (2 * (1 - c))) := by
  obtain ⟨hx0, hx1⟩ := hx
  have hc1 : c < 1 := by linarith
  have hH := hiding_bound hProt hiid hw hw1 (hC 0) hChide hCc hc0 hc1 hcert
  have hw0 : 0 ≤ w 0 := hw 0
  have hw01 : w 0 ≤ 1 := weight_zero_le_one hw hw1
  have hx2p : (0:ℝ) < x^2 := by positivity
  have hx2l : x^2 < 1 := by nlinarith
  have hB0 : 0 < h (x^2) := h_pos hx2p hx2l
  have hB1 : h (x^2) ≤ 1 := h_le_one _
  have hcond : h x ≤ w 0 * h (x^2) + 1 - w 0 := by
    linarith [mul_nonneg (sub_nonneg.2 hw01) (sub_nonneg.2 hB1)]
  have hD := diagonal_bound hProt hiid hw hw1 hC hCc hc1 hcert ⟨hx0, le_of_lt hx1⟩ hcond
  have hden : 0 < 2*(1-c) := by linarith
  have hDpos : 0 < w 0 * h (x^2) + 1 - w 0 := lt_of_lt_of_le (h_pos hx0 hx1) hcond
  have hD0pos : 0 < 1 - (1 - h (x^2)) / (2*(1-c)) := by
    rw [sub_pos, div_lt_one hden]; linarith
  have hw0lb : 1/(2*(1-c)) ≤ w 0 := by rw [div_le_iff₀ hden]; linarith
  have hDD0 : w 0 * h (x^2) + 1 - w 0 ≤ 1 - (1 - h (x^2)) / (2*(1-c)) := by
    have hkey : (1 - h (x^2))/(2*(1-c)) ≤ w 0 * (1 - h (x^2)) := by
      rw [div_le_iff₀ hden]
      linarith [mul_nonneg (sub_nonneg.2 hB1)
        (by linarith [hH] : (0:ℝ) ≤ 2*w 0*(1-c) - 1)]
    linarith
  have hnum : 0 ≤ x * h x := mul_nonneg (le_of_lt hx0) (h_nonneg (le_of_lt hx0) (le_of_lt hx1))
  have hinv : 1/(1 - (1 - h (x^2)) / (2*(1-c))) ≤ 1/(w 0 * h (x^2) + 1 - w 0) :=
    one_div_le_one_div_of_le hDpos hDD0
  have hcmp : x * h x / (1 - (1 - h (x^2)) / (2*(1-c)))
      ≤ x * h x / (w 0 * h (x^2) + 1 - w 0) :=
    calc x * h x / (1 - (1 - h (x^2)) / (2*(1-c)))
        = x * h x * (1/(1 - (1 - h (x^2)) / (2*(1-c)))) := by ring
      _ ≤ x * h x * (1/(w 0 * h (x^2) + 1 - w 0)) := mul_le_mul_of_nonneg_left hinv hnum
      _ = x * h x / (w 0 * h (x^2) + 1 - w 0) := by ring
  linarith [hD, hcmp]

/-- **The numerical corollary of Theorem 3.4.**  Under the hypotheses of the section
(and `c ≤ 1/2`, as in Proposition 2.2), the certified constant satisfies `c ≤ 0.3829`.

This lies strictly below the product ceiling `c_ceil = 0.3830993…` of Theorem 3.1; the
exact value in the paper is `c** = 0.382885260…`.  The adversary is the two-point law at
`x = 0.6909`. -/
theorem refined_ceiling_numeric
    (hProt : ∀ k, IsProtocol (Prot k)) (hiid : ∀ x y, Prot 0 x y = x * y)
    (hw : ∀ k, 0 ≤ w k) (hw1 : ∑ k, w k = 1) (hC : ∀ k, ContainsProduct (C k))
    (hChide : ∀ k, k ≠ 0 → AdmitsHiding (C k))
    (hCc : 1 ≤ Cc) (hc0 : 0 < c) (hchalf : c ≤ 1/2) (hcert : Certifies Prot w C Cc c) :
    c ≤ 3829/10000 := by
  by_contra hcon
  push_neg at hcon
  have hxm : ((6909:ℝ)/10000) ∈ Set.Ioo (0:ℝ) 1 := ⟨by norm_num, by norm_num⟩
  have hBl : (9985182/10000000:ℝ) ≤ h (((6909:ℝ)/10000)^2) := by rw [xsq_eq]; exact hx2_lower
  have hBu : h (((6909:ℝ)/10000)^2) ≤ (998519/1000000:ℝ) := by rw [xsq_eq]; exact hx2_upper
  have hmono : h ((6909:ℝ)/10000) ≤ h (((6909:ℝ)/10000)^2) := by
    linarith [hx_upper, hBl]
  have hfp := fixed_point_form hProt hiid hw hw1 hC hChide hCc hc0 hchalf hcert hxm hmono
  have hc1 : c < 1 := by linarith
  have hden : 0 < 2*(1-c) := by linarith
  have hB1 : h (((6909:ℝ)/10000)^2) ≤ 1 := h_le_one _
  have hD0pos : 0 < 1 - (1 - h (((6909:ℝ)/10000)^2)) / (2*(1-c)) := by
    rw [sub_pos, div_lt_one hden]; linarith
  have hne : (1:ℝ) - c ≠ 0 := by linarith
  have key : ((6909:ℝ)/10000) * h ((6909:ℝ)/10000)
      ≤ (1 - c) - (1 - h (((6909:ℝ)/10000)^2))/2 := by
    have e1 : ((6909:ℝ)/10000) * h ((6909:ℝ)/10000)
        / (1 - (1 - h (((6909:ℝ)/10000)^2)) / (2*(1-c))) ≤ 1 - c := by linarith [hfp]
    rw [div_le_iff₀ hD0pos] at e1
    have e2 : (1 - c) * (1 - (1 - h (((6909:ℝ)/10000)^2)) / (2*(1-c)))
        = (1 - c) - (1 - h (((6909:ℝ)/10000)^2))/2 := by
      field_simp
    linarith [e1, e2.le, e2.ge]
  linarith [key, hx_lower, hBu, hcon]

end UnionClosedCeiling
\end{Verbatim}

\subsection{\texttt{UnionClosedCeiling/Classes.lean}}
\begin{Verbatim}[fontsize=\scriptsize,formatcom=\leanmono,samepage=false]
import UnionClosedCeiling.Framework

/-!
# Lemma 3.3 — the classes

> **Lemma 3.3.**  The classes of the i.i.d. protocol, of Sawin's protocol, of Yu's
> protocols for every `ρ > 0`, and of every conditionally-i.i.d. protocol contain the
> product law.  All but the first admit hiding.

Formalised here, in the finite model:

* the singleton class `{μ ⊗ μ}` of the i.i.d. protocol (`iidClass`) contains the product;
* the class of **all couplings** (`allCouplings`, Sawin's class) contains the product and
  admits hiding;
* the class of **mixtures of products** (`mixtureOfProducts`, the class of every
  conditionally-i.i.d. protocol) contains the product and admits hiding.

The witness used for hiding is the paper's mixture-of-two-products coupling
`P = q δ₁⊗δ₁ + (1-q) P₀⊗P₀`, `P₀ = (1-δ)δ₀ + δ δ_y` (`hidingCoupling`).  For the class of
all couplings the paper offers the maximal-correlation witness instead; since that class
imposes no constraint, the same mixture witness serves, and we use it for both.

**Not formalised:** the maximal-correlation classes of Yu's protocols.  Their hiding
witness requires the continuity of singular values in the matrix entries, which is a
genuinely different argument; see the README.
-/

namespace UnionClosedCeiling

open Finset

/-- The class of *all* couplings (Sawin's class). -/
def allCouplings : ∀ n : ℕ, (Fin n → ℝ) → (Fin n → ℝ) → (Fin n → Fin n → ℝ) → Prop :=
  fun _ _ _ _ => True

/-- The singleton class `{μ ⊗ μ}` of the i.i.d. protocol. -/
def iidClass : ∀ n : ℕ, (Fin n → ℝ) → (Fin n → ℝ) → (Fin n → Fin n → ℝ) → Prop :=
  fun _ _ m P => P = productCoupling m

/-- The class of mixtures of product laws: `P = ∑ⱼ πⱼ νⱼ ⊗ νⱼ` with `∑ⱼ πⱼ νⱼ = μ`.
This is the class `C₃(μ)` of the conditionally-i.i.d. protocols. -/
def mixtureOfProducts : ∀ n : ℕ, (Fin n → ℝ) → (Fin n → ℝ) → (Fin n → Fin n → ℝ) → Prop :=
  fun n _ m P => ∃ (J : ℕ) (π : Fin J → ℝ) (ν : Fin J → Fin n → ℝ),
    (∀ j, 0 ≤ π j) ∧ (∑ j, π j = 1) ∧
    (∀ j i, 0 ≤ ν j i) ∧ (∀ j, ∑ i, ν j i = 1) ∧
    (∀ i, ∑ j, π j * ν j i = m i) ∧
    (∀ i i', P i i' = ∑ j, π j * (ν j i * ν j i'))

/-! ### Products -/

lemma containsProduct_allCouplings : ContainsProduct allCouplings :=
  fun _ _ _ _ => trivial

lemma containsProduct_iidClass : ContainsProduct iidClass :=
  fun _ _ _ _ => rfl

lemma containsProduct_mixtureOfProducts : ContainsProduct mixtureOfProducts := by
  intro n a m hlaw
  refine ⟨1, ![1], ![m], ?_, ?_, ?_, ?_, ?_, ?_⟩
  · intro j; fin_cases j; norm_num
  · simp
  · intro j i; fin_cases j; exact hlaw.2.1 i
  · intro j; fin_cases j; exact hlaw.2.2
  · intro i; simp
  · intro i i'; simp [productCoupling]

/-! ### Hiding -/

/-- The paper's hiding coupling of `μ_{q,y,δ}` with itself:
`q δ₁ ⊗ δ₁ + (1-q) P₀ ⊗ P₀` with `P₀ = (1-δ) δ₀ + δ δ_y`.
In our indexing (`0 ↦ 1`, `1 ↦ 0`, `2 ↦ y`), `δ₁ = ![1,0,0]` and `P₀ = ![0,1-δ,δ]`. -/
noncomputable def hidingComp1 : Fin 3 → ℝ := ![1, 0, 0]

/-- The second component `P₀ = (1-δ) δ₀ + δ δ_y` of the hiding coupling. -/
noncomputable def hidingComp2 (δ : ℝ) : Fin 3 → ℝ := ![0, 1 - δ, δ]

/-- The hiding coupling `P = q δ₁⊗δ₁ + (1-q) P₀⊗P₀`. -/
noncomputable def hidingCoupling (q δ : ℝ) : Fin 3 → Fin 3 → ℝ :=
  fun i j => q * (hidingComp1 i * hidingComp1 j) + (1-q) * (hidingComp2 δ i * hidingComp2 δ j)

lemma hidingCoupling_isCoupling {q δ : ℝ} (hq0 : 0 ≤ q) (hq1 : q ≤ 1)
    (hδ0 : 0 ≤ δ) (hδ1 : δ ≤ 1) : IsCoupling (hidingWeights q δ) (hidingCoupling q δ) := by
  have hc1 : ∀ i, 0 ≤ hidingComp1 i := by
    intro i; fin_cases i <;> simp [hidingComp1]
  have hc2 : ∀ i, 0 ≤ hidingComp2 δ i := by
    intro i; fin_cases i <;> simp [hidingComp2] <;> linarith
  refine ⟨?_, ?_, ?_⟩
  · intro i j
    have := mul_nonneg (hc1 i) (hc1 j)
    have := mul_nonneg (hc2 i) (hc2 j)
    unfold hidingCoupling
    have h1 : 0 ≤ q * (hidingComp1 i * hidingComp1 j) :=
      mul_nonneg hq0 (mul_nonneg (hc1 i) (hc1 j))
    have h2 : 0 ≤ (1-q) * (hidingComp2 δ i * hidingComp2 δ j) :=
      mul_nonneg (by linarith) (mul_nonneg (hc2 i) (hc2 j))
    linarith
  · intro i
    fin_cases i <;>
      simp [hidingCoupling, hidingComp1, hidingComp2, hidingWeights, Fin.sum_univ_three] <;> ring
  · intro j
    fin_cases j <;>
      simp [hidingCoupling, hidingComp1, hidingComp2, hidingWeights, Fin.sum_univ_three] <;> ring

lemma hidingCoupling_mixture {q δ : ℝ} (hq0 : 0 ≤ q) (hq1 : q ≤ 1)
    (hδ0 : 0 ≤ δ) (hδ1 : δ ≤ 1) (y : ℝ) :
    mixtureOfProducts 3 (hidingAtoms y) (hidingWeights q δ) (hidingCoupling q δ) := by
  refine ⟨2, ![q, 1-q], ![hidingComp1, hidingComp2 δ], ?_, ?_, ?_, ?_, ?_, ?_⟩
  · intro j; fin_cases j
    · exact hq0
    · show (0:ℝ) ≤ 1 - q
      linarith
  · simp [Fin.sum_univ_two]
  · intro j i; fin_cases j <;> fin_cases i <;>
      simp [hidingComp1, hidingComp2] <;> linarith
  · intro j; fin_cases j <;> simp [hidingComp1, hidingComp2, Fin.sum_univ_three]
  · intro i; fin_cases i <;>
      simp [hidingComp1, hidingComp2, hidingWeights, Fin.sum_univ_two]
  · intro i i'; simp [hidingCoupling, Fin.sum_univ_two]

lemma hidingCoupling_zeros {q δ : ℝ} : hidingCoupling q δ 2 0 = 0 ∧ hidingCoupling q δ 0 2 = 0 := by
  constructor <;> simp [hidingCoupling, hidingComp1, hidingComp2]

lemma hidingCoupling_diag {q δ : ℝ} (hq0 : 0 ≤ q) (hδ0 : 0 ≤ δ) :
    hidingCoupling q δ 2 2 ≤ δ^2 := by
  have : hidingCoupling q δ 2 2 = (1-q) * (δ * δ) := by
    simp [hidingCoupling, hidingComp1, hidingComp2]
  rw [this]
  nlinarith [mul_nonneg hδ0 hδ0]

/-- The class of all couplings admits hiding. -/
lemma admitsHiding_allCouplings : AdmitsHiding allCouplings := by
  intro q y hq hy
  refine ⟨1, one_pos, fun δ hδ0 hδ1 => ⟨hidingCoupling q δ, ?_, trivial, ?_, ?_, ?_⟩⟩
  · exact hidingCoupling_isCoupling (le_of_lt hq.1) (le_of_lt hq.2) (le_of_lt hδ0)
      (le_of_lt hδ1)
  · exact hidingCoupling_zeros.1
  · exact hidingCoupling_zeros.2
  · exact hidingCoupling_diag (le_of_lt hq.1) (le_of_lt hδ0)

/-- The class of mixtures of products admits hiding. -/
lemma admitsHiding_mixtureOfProducts : AdmitsHiding mixtureOfProducts := by
  intro q y hq hy
  refine ⟨1, one_pos, fun δ hδ0 hδ1 => ⟨hidingCoupling q δ, ?_, ?_, ?_, ?_, ?_⟩⟩
  · exact hidingCoupling_isCoupling (le_of_lt hq.1) (le_of_lt hq.2) (le_of_lt hδ0)
      (le_of_lt hδ1)
  · exact hidingCoupling_mixture (le_of_lt hq.1) (le_of_lt hq.2) (le_of_lt hδ0)
      (le_of_lt hδ1) y
  · exact hidingCoupling_zeros.1
  · exact hidingCoupling_zeros.2
  · exact hidingCoupling_diag (le_of_lt hq.1) (le_of_lt hδ0)

end UnionClosedCeiling
\end{Verbatim}

\subsection{\texttt{UnionClosedCeiling.lean}}
\begin{Verbatim}[fontsize=\scriptsize,formatcom=\leanmono,samepage=false]
import UnionClosedCeiling.Entropy
import UnionClosedCeiling.Framework
import UnionClosedCeiling.Ceiling
import UnionClosedCeiling.Refined
import UnionClosedCeiling.Classes

/-!
# A ceiling for single-letter protocol arguments (Sections 2.2 and 3 of the paper)

Main results:

* `UnionClosedCeiling.product_ceiling`   -- Theorem 3.1 (product ceiling)
* `UnionClosedCeiling.hiding_bound`      -- Theorem 3.4, inequality (H)
* `UnionClosedCeiling.diagonal_bound`    -- Theorem 3.4, inequality (D)
* `UnionClosedCeiling.fixed_point_form`  -- Theorem 3.4, the fixed-point form
* `UnionClosedCeiling.refined_ceiling_numeric` -- Theorem 3.4, `c ≤ 0.3829`
* `UnionClosedCeiling.containsProduct_*`, `UnionClosedCeiling.admitsHiding_*` -- Lemma 3.3
-/

#print axioms UnionClosedCeiling.product_ceiling
#print axioms UnionClosedCeiling.hiding_bound
#print axioms UnionClosedCeiling.diagonal_bound
#print axioms UnionClosedCeiling.fixed_point_form
#print axioms UnionClosedCeiling.refined_ceiling_numeric
#print axioms UnionClosedCeiling.containsProduct_iidClass
#print axioms UnionClosedCeiling.containsProduct_allCouplings
#print axioms UnionClosedCeiling.containsProduct_mixtureOfProducts
#print axioms UnionClosedCeiling.admitsHiding_allCouplings
#print axioms UnionClosedCeiling.admitsHiding_mixtureOfProducts
\end{Verbatim}

\subsection{\texttt{AxiomCheck.lean}}
\begin{Verbatim}[fontsize=\scriptsize,formatcom=\leanmono,samepage=false]
import UnionClosedCeiling

#print axioms UnionClosedCeiling.product_ceiling
#print axioms UnionClosedCeiling.hiding_bound
#print axioms UnionClosedCeiling.diagonal_bound
#print axioms UnionClosedCeiling.fixed_point_form
#print axioms UnionClosedCeiling.refined_ceiling_numeric
#print axioms UnionClosedCeiling.containsProduct_iidClass
#print axioms UnionClosedCeiling.containsProduct_allCouplings
#print axioms UnionClosedCeiling.containsProduct_mixtureOfProducts
#print axioms UnionClosedCeiling.admitsHiding_allCouplings
#print axioms UnionClosedCeiling.admitsHiding_mixtureOfProducts
\end{Verbatim}

\fussy

\begin{thebibliography}{99}
\bibitem{AHS} R.~Alweiss, B.~Huang and M.~Sellke, Improved lower bound for the union-closed sets conjecture, arXiv:2211.11731 (2022).
\bibitem{BruhnSchaudt} H.~Bruhn and O.~Schaudt, The journey of the union-closed sets conjecture, \emph{Graphs Combin.} 31 (2015), 2043--2074.
\bibitem{Cambie} S.~Cambie, Better bounds for the union-closed sets conjecture using the entropy approach, arXiv:2212.12500 (2022; revised 2025).
\bibitem{CL} Z.~Chase and S.~Lovett, Approximate union closed conjecture, arXiv:2211.11689 (2022).
\bibitem{Frankl} P.~Frankl, Extremal set systems, in: \emph{Handbook of Combinatorics}, Vol.~2, Elsevier, Amsterdam, 1995, pp.~1293--1329.
\bibitem{Gilmer} J.~Gilmer, A constant lower bound for the union-closed sets conjecture, arXiv:2211.09055 (2022).
\bibitem{Karr} A.~F.~Karr, Extreme points of certain sets of probability measures, with applications, \emph{Math. Oper. Res.} 8 (1983), 74--85.
\bibitem{Liu} J.~Liu, Improving the lower bound for the union-closed sets conjecture via conditionally IID coupling, in: \emph{2024 IEEE International Symposium on Information Theory (ISIT)}; arXiv:2306.08824 (2023).
\bibitem{Pebody} L.~Pebody, Extension of a method of Gilmer, arXiv:2211.13139 (2022).
\bibitem{Sawin} W.~Sawin, An improved lower bound for the union-closed set conjecture, arXiv:2211.11504 (2022).
\bibitem{Winkler} G.~Winkler, Extreme points of moment sets, \emph{Math. Oper. Res.} 13 (1988), 581--587.
\bibitem{Yu} L.~Yu, Dimension-free bounds for the union-closed sets conjecture, \emph{Entropy} 25 (2023), no.~5, 767; arXiv:2212.00658.
\end{thebibliography}
\end{document}
